\section{Supplement on homological algebra}
\label{section:0.11}

\subsection{Review of spectral sequences}
\label{subsection:0.11.1}

\begin{env}[11.1.1]
\label{0.11.1.1}
In what follows we will use a notion of spectral sequence more general than
the one defined in (T, 2.4); keeping the notation of (T, 2.4), we will
call a \emph{spectral sequence} in an abelian category $\mathcal C$ a
system E consisting of the following elements:
\begin{enumerate}
  \item[(a)] A family $(E_r^{pq})$ of objects of $\mathcal C$ defined for
    $p\in\mathbb Z$, $q\in\mathbb Z$ and $r\geq2$.
  \item[(b)] A family of morphisms
    $d_r^{pq}:E_r^{pq}\to E_r^{p+r,q-r+1}$ such that
    $d_r^{p+r,q-r+1}d_r^{pq}=0$. We set
    $Z_{r+1}(E_r^{pq})=\operatorname{Ker}(d_r^{pq})$,
    $B_{r+1}(E_r^{pq})=\operatorname{Im}(d_r^{p-r,q+r-1})$, so that
    \[
      B_{r+1}(E_r^{pq})\subset Z_{r+1}(E_r^{pq})\subset E_r^{pq}.
    \]
  \item[(c)] A family of isomorphisms
    $\alpha_r^{pq}:Z_{r+1}(E_r^{pq})/B_{r+1}(E_r^{pq})
    \xrightarrow{\sim}E_{r+1}^{pq}$.
\end{enumerate}

We then define for $k\geq r+1$, by induction on $k$, the subobjects
$B_k(E_r^{pq})$ and $Z_k(E_r^{pq})$ as the inverse images, under the
canonical morphism $E_r^{pq}\to E_r^{pq}/B_{r+1}(E_r^{pq})$, of the subobjects of
this quotient identified by $\alpha_r^{pq}$ with the subobjects $B_k(E_{r+1}^{pq})$ and
$Z_k(E_{r+1}^{pq})$ respectively. It is clear that we then have, up to
isomorphism,
\[
\label{0.11.1.1.1}
  Z_k(E_r^{pq})/B_k(E_r^{pq})=E_k^{pq}
  \qquad\text{for }k\geq r+1,
  \tag{11.1.1.1}
\]
and, if we further set $B_r(E_r^{pq})=0$ and $Z_r(E_r^{pq})=E_r^{pq}$, we have
the inclusion relations
\[
\label{0.11.1.1.2}
  0=B_r(E_r^{pq})\subset B_{r+1}(E_r^{pq})\subset B_{r+2}(E_r^{pq})
  \subset\cdots\subset Z_{r+2}(E_r^{pq})\subset Z_{r+1}(E_r^{pq})
  \subset Z_r(E_r^{pq})=E_r^{pq}.
  \tag{11.1.1.2}
\]

The other elements of the data of E are then:
\begin{enumerate}
  \item[(d)] Two subobjects $B_\infty(E_2^{pq})$ and
    $Z_\infty(E_2^{pq})$ of $E_2^{pq}$ such that
    $B_\infty(E_2^{pq})\subset Z_\infty(E_2^{pq})$ and, for every $k\geq2$,
    \[
      B_k(E_2^{pq})\subset B_\infty(E_2^{pq})
      \qquad\text{and}\qquad
      Z_\infty(E_2^{pq})\subset Z_k(E_2^{pq}).
    \]
    We set
    \[
    \label{0.11.1.1.3}
      E_\infty^{pq}=Z_\infty(E_2^{pq})/B_\infty(E_2^{pq}).
      \tag{11.1.1.3}
    \]
  \item[(e)]
    \oldpage[0\textsubscript{III}]{24}
    A family $(E^n)$ of objects of $\mathcal C$, each of which is equipped with a
    \emph{decreasing filtration} $(F^p(E^n))_{p\in\mathbb Z}$. As usual we
    denote by $\operatorname{gr}(E^n)$ the graded object associated to the
    filtered object $E^n$, the direct sum of the
    $\operatorname{gr}_p(E^n)=F^p(E^n)/F^{p+1}(E^n)$.
  \item[(f)] For every pair $(p,q)\in\mathbb Z\times\mathbb Z$, an
    isomorphism
    $\beta^{pq}:E_\infty^{pq}\xrightarrow{\sim}
    \operatorname{gr}_p(E^{p+q})$.
\end{enumerate}

The family $(E^n)$, without the filtrations, is called \emph{the abutment}
of the spectral sequence E.

Suppose that the category $\mathcal C$ admits infinite direct sums,
or that for every $r\geq2$ and every $n\in\mathbb Z$, there are only finitely
many pairs $(p,q)$ such that $p+q=n$ and $E_r^{pq}\ne0$ (it suffices that this be
so for $r=2$). Then the
\[
  E_r^{(n)}=\sum_{p+q=n}E_r^{pq}
\]
are defined, and, denoting by $d_r^{(n)}$ the morphism
$E_r^{(n)}\to E_r^{(n+1)}$ whose restriction to $E_r^{pq}$ is $d_r^{pq}$
for each pair $(p,q)$ such that $p+q=n$, we have
$d_r^{(n+1)}\circ d_r^{(n)}=0$, in other words $(E_r^{(n)})_{n\in\mathbb Z}$
is a \emph{complex $E_r^{(\bullet)}$} in $\mathcal C$, with differential of
degree $+1$, and it follows from c) that
$\mathrm H^n(E_r^{(\bullet)})$ is isomorphic to $E_{r+1}^{(n)}$ for every
$r\geq2$.
\end{env}

\begin{env}[11.1.2]
\label{0.11.1.2}
A \emph{morphism $u:E\to E'$} from a spectral sequence E to a spectral
sequence $E'=(E_r^{\prime pq},E^{\prime n})$ consists of systems of
morphisms $u_r^{pq}:E_r^{pq}\to E_r^{\prime pq}$,
$u^n:E^n\to E^{\prime n}$, the $u^n$ being compatible with the filtrations
of $E^n$ and $E^{\prime n}$, and the diagrams
\[
  \xymatrix{
    E_r^{pq}\ar[r]^-{d_r^{pq}}\ar[d]_{u_r^{pq}} &
    E_r^{p+r,q-r+1}\ar[d]^{u_r^{p+r,q-r+1}}\\
    E_r^{\prime pq}\ar[r]_{d_r^{\prime pq}} &
    E_r^{\prime p+r,q-r+1}
  }
\]
being commutative; moreover, by passing to quotients, $u_r^{pq}$ gives a
morphism
\[
  \overline u_r^{pq}:
  Z_{r+1}(E_r^{pq})/B_{r+1}(E_r^{pq})\longrightarrow
  Z_{r+1}(E_r^{\prime pq})/B_{r+1}(E_r^{\prime pq})
\]
and we must have
$\alpha_r^{\prime pq}\circ\overline u_r^{pq}
=u_{r+1}^{pq}\circ\alpha_r^{pq}$; finally, we must have
$u_2^{pq}(B_\infty(E_2^{pq}))\subset B_\infty(E_2^{\prime pq})$,
$u_2^{pq}(Z_\infty(E_2^{pq}))\subset Z_\infty(E_2^{\prime pq})$; by
passing to quotients, $u_2^{pq}$ then gives a morphism
$u_\infty^{pq}:E_\infty^{pq}\to E_\infty^{\prime pq}$, and the diagram
\[
  \xymatrix{
    E_\infty^{pq}\ar[r]^-{u_\infty^{pq}}\ar[d]_{\beta^{pq}} &
    E_\infty^{\prime pq}\ar[d]^{\beta^{\prime pq}}\\
    \operatorname{gr}_p(E^{p+q})
      \ar[r]_{\operatorname{gr}_p(u^{p+q})} &
    \operatorname{gr}_p(E^{\prime p+q})
  }
\]
must be commutative.

The preceding definitions show, by induction on $r$, that if the
$u_2^{pq}$ are \emph{isomorphisms}, then so are the $u_r^{pq}$
for $r\geq2$; \emph{if we know moreover that}
$u_2^{pq}(B_\infty(E_2^{pq}))=B_\infty(E_2^{\prime pq})$ and
$u_2^{pq}(Z_\infty(E_2^{pq}))=Z_\infty(E_2^{\prime pq})$ \emph{and that the
$u^n$ are isomorphisms}, then we can conclude that $u$ is an
\emph{isomorphism}.
\end{env}

\begin{env}[11.1.3]
\label{0.11.1.3}
\oldpage[0\textsubscript{III}]{25}
Recall that if $(F^p(X))_{p\in\mathbb Z}$ is a (decreasing) filtration
of an object $X\in\mathcal C$, we say that this filtration is \emph{separated}
if $\inf(F^p(X))=0$, \emph{discrete} if there exists a $p$ such that $F^p(X)=0$,
\emph{exhaustive} (or \emph{co-separated}) if $\sup(F^p(X))=X$, and
\emph{co-discrete} if there exists $p$ such that $F^p(X)=X$.

We will say that a spectral sequence $E=(E_r^{pq},E^n)$ is \emph{weakly
convergent} if we have
\[
  B_\infty(E_2^{pq})=\sup_k(B_k(E_2^{pq})),\qquad
  Z_\infty(E_2^{pq})=\inf_k(Z_k(E_2^{pq}))
\]
(in other words the objects $B_\infty(E_2^{pq})$ and $Z_\infty(E_2^{pq})$
are determined by the data a) to c) of the spectral sequence E). We will say
that the spectral sequence E is \emph{regular} if it is weakly
convergent and if moreover:
\begin{enumerate}
  \item[$1^\circ$] For every pair $(p,q)$, the decreasing sequence
    $(Z_k(E_2^{pq}))_{k\geq2}$ is \emph{stationary}; the hypothesis that E
    is weakly convergent then implies
    $Z_\infty(E_2^{pq})=Z_k(E_2^{pq})$ for $k$ sufficiently large (depending on
    $p$ and $q$).
  \item[$2^\circ$] For every $n$, the filtration
    $(F^p(E^n))_{p\in\mathbb Z}$ of $E^n$ is \emph{discrete} and
    \emph{exhaustive}.
\end{enumerate}

We say that the spectral sequence E is \emph{coregular} if it is
weakly convergent and if moreover:
\begin{enumerate}
  \item[$3^\circ$] For every pair $(p,q)$, the increasing sequence
    $(B_k(E_2^{pq}))_{k\geq2}$ is \emph{stationary}, which implies
    $B_\infty(E_2^{pq})=B_k(E_2^{pq})$, and consequently
    $E_\infty^{pq}=\inf E_k^{pq}$.
  \item[$4^\circ$] For every $n$, the filtration of $E^n$ is
    \emph{co-discrete}.
\end{enumerate}

Finally, we say that E is \emph{biregular} if it is both regular
and coregular, in other words if the following conditions hold:
\begin{enumerate}
  \item[(a)] For every pair $(p,q)$, the sequences
    $(B_k(E_2^{pq}))_{k\geq2}$ and $(Z_k(E_2^{pq}))_{k\geq2}$ are
    \emph{stationary} and we have
    $B_\infty(E_2^{pq})=B_k(E_2^{pq})$ and
    $Z_\infty(E_2^{pq})=Z_k(E_2^{pq})$ for $k$ sufficiently large (which implies
    $E_\infty^{pq}=E_k^{pq}$).
  \item[(b)] For every $n$, the filtration $(F^p(E^n))_{p\in\mathbb Z}$ is
    \emph{discrete} and \emph{co-discrete} (which is also expressed by saying
    that it is \emph{finite}).
\end{enumerate}

The spectral sequences defined in (T, 2.4) are thus the biregular
spectral sequences.
\end{env}

\begin{env}[11.1.4]
\label{0.11.1.4}
Suppose that in the category $\mathcal C$ filtered inductive limits
exist and that the functor $\varinjlim$ is \emph{exact} (which amounts to
saying that axiom AB 5) of (T, 1.5) holds (cf. T, 1.8)). The condition
that the filtration $(F^p(X))_{p\in\mathbb Z}$ of an object $X\in\mathcal C$ be
exhaustive can then also be written
$\varinjlim_{p\to-\infty}F^p(X)=X$. If a spectral sequence E is weakly
convergent, we have
$B_\infty(E_2^{pq})=\varinjlim_{k\to\infty}B_k(E_2^{pq})$; if moreover
$u:E\to E'$ is a morphism from E to a weakly convergent spectral sequence $E'$ of
$\mathcal C$, we have
$u_2^{pq}(B_\infty(E_2^{pq}))=B_\infty(E_2^{\prime pq})$ because of
the exactness of $\varinjlim$. Moreover:
\end{env}

\begin{proposition}[11.1.5]
\label{0.11.1.5}
Let $\mathcal C$ be an abelian category in which filtered inductive limits
are exact, E, $E'$ two regular spectral sequences
of $\mathcal C$, and $u:E\to E'$ a morphism of spectral sequences. If the
$u_2^{pq}$ are isomorphisms, then so is $u$.
\end{proposition}

\begin{proof}
We already know (\hyperref[0.11.1.2]{11.1.2}) that the $u_r^{pq}$ are
isomorphisms and that
\[
  u_2^{pq}(B_\infty(E_2^{pq}))=B_\infty(E_2^{\prime pq})~;
\]
\oldpage[0\textsubscript{III}]{26}
the hypothesis that E and $E'$ are regular also implies
$u_2^{pq}(Z_\infty(E_2^{pq}))=Z_\infty(E_2^{\prime pq})$, and since
$u_2^{pq}$ is an isomorphism, so is $u_\infty^{pq}$; we therefore
conclude that $\operatorname{gr}_p(u^{p+q})$ is also an isomorphism.
But since the filtrations of the $E^n$ and of the $E^{\prime n}$ are discrete and
exhaustive, this implies that the $u^n$ are also isomorphisms
(Bourbaki, \emph{Alg. comm.}, chap. III, \S~2, n\textsuperscript{o} 8,
th. 1).
\end{proof}

\begin{env}[11.1.6]
\label{0.11.1.6}
It follows from (\hyperref[0.11.1.1.2]{11.1.1.2}) and from the definition
(\hyperref[0.11.1.1.3]{11.1.1.3}) that if, in a spectral sequence E, we have
$E_r^{pq}=0$, then $E_k^{pq}=0$ for $k\geq r$ and $E_\infty^{pq}=0$. We say
that a spectral sequence is \emph{degenerate} if there exist an integer $r\geq2$ and,
for every integer $n\in\mathbb Z$, an integer $q(n)$ such that
$E_r^{n-q,q}=0$ for every $q\ne q(n)$. We first deduce from the preceding
remark that we also have $E_k^{n-q,q}=0$ for $k\geq r$ (including
$k=\infty$) and $q\ne q(n)$. Moreover, the definition of $E_{r+1}^{pq}$ shows
that we have
$E_{r+1}^{n-q(n),q(n)}=E_r^{n-q(n),q(n)}$; if E is weakly convergent,
we therefore also have
$E_\infty^{n-q(n),q(n)}=E_r^{n-q(n),q(n)}$; in other words, for every
$n\in\mathbb Z$, $\operatorname{gr}_p(E^n)=0$ for $p\ne n-q(n)$ and
$\operatorname{gr}_{n-q(n)}(E^n)=E_r^{n-q(n),q(n)}$. If moreover the filtration
of $E^n$ is discrete and exhaustive, the sequence E is regular and we have
$E^n=E_r^{n-q(n),q(n)}$ up to isomorphism.
\end{env}

\begin{env}[11.1.7]
\label{0.11.1.7}
Suppose that in the category $\mathcal C$ filtered inductive limits
exist and are exact, and let $(E_\lambda,u_{\mu\lambda})$ be an
inductive system (over a filtered index set) of spectral sequences of
$\mathcal C$. Then the \emph{inductive limit} of this inductive system exists
in the additive category of spectral sequences of objects of $\mathcal C$:
to see this, it suffices to define $E_r^{pq}$, $d_r^{pq}$, $\alpha_r^{pq}$,
$B_\infty(E_2^{pq})$, $Z_\infty(E_2^{pq})$, $E^n$, $F^p(E^n)$ and
$\beta^{pq}$ as the respective inductive limits of $E_{r,\lambda}^{pq}$,
$d_{r,\lambda}^{pq}$, $\alpha_{r,\lambda}^{pq}$,
$B_\infty(E_{2,\lambda}^{pq})$, $Z_\infty(E_{2,\lambda}^{pq})$,
$E_\lambda^n$, $F^p(E_\lambda^n)$ and $\beta_\lambda^{pq}$; the verification
of the conditions of (\hyperref[0.11.1.1]{11.1.1}) follows from the exactness of the
functor $\varinjlim$ in $\mathcal C$.
\end{env}

\begin{remark}[11.1.8]
\label{0.11.1.8}
Suppose that the category $\mathcal C$ is the category of $A$-modules over
a \emph{Noetherian} ring $A$ (resp. over a ring $A$). Then the
definitions of (\hyperref[0.11.1.1]{11.1.1}) show that if, for a given $r$,
the $E_r^{pq}$ are $A$-modules \emph{of finite type} (resp. of
\emph{finite length}), then so are all the modules $E_s^{pq}$ for
$s\geq r$, as well as the $E_\infty^{pq}$. If moreover the filtration of
the abutment $(E^n)$ is discrete and co-discrete for every $n$, we
conclude that each of the $E^n$ is also an $A$-module \emph{of finite type}
(resp. of \emph{finite length}).
\end{remark}

\begin{env}[11.1.9]
\label{0.11.1.9}
We will have to consider conditions ensuring that a spectral sequence E is
biregular in a ``uniform'' way in $p+q=n$. We will then use the following
lemma:
\end{env}

\begin{lemma}[11.1.10]
\label{0.11.1.10}
Let $(E_r^{pq})$ be a family of objects of $\mathcal C$ related by the data
a), b), c) of (\hyperref[0.11.1.1]{11.1.1}). For a fixed integer $n$, the
following properties are equivalent:
\begin{enumerate}
  \item[(a)] There exists an integer $r(n)$ such that for $r\geq r(n)$,
    $p+q=n$ or $p+q=n-1$, the morphisms $d_r^{pq}$ are all zero.
  \item[(b)] There exists an integer $r(n)$ such that for $p+q=n$ or $p+q=n+1$,
    we have $B_r(E_2^{pq})=B_s(E_2^{pq})$ for $s\geq r\geq r(n)$.
  \item[(c)] There exists an integer $r(n)$ such that for $p+q=n$ or $p+q=n-1$,
    we have $Z_r(E_2^{pq})=Z_s(E_2^{pq})$ for $s\geq r\geq r(n)$.
  \item[(d)] There exists an integer $r(n)$ such that for $p+q=n$, we have
    $B_r(E_2^{pq})=B_s(E_2^{pq})$ and $Z_r(E_2^{pq})=Z_s(E_2^{pq})$ for
    $s\geq r\geq r(n)$.
\end{enumerate}
\end{lemma}

\begin{proof}
\oldpage[0\textsubscript{III}]{27}
Indeed, by conditions a), b), c) of
(\hyperref[0.11.1.1]{11.1.1}), saying that
$Z_{r+1}(E_2^{pq})=Z_r(E_2^{pq})$ is equivalent to saying that $d_r^{pq}=0$, and saying
that
$B_r(E_2^{p+r,q-r+1})=B_{r+1}(E_2^{p+r,q-r+1})$ is also equivalent to saying that
$d_r^{pq}=0$; the lemma follows immediately from this remark.
\end{proof}
\subsection{The spectral sequence of a filtered complex}
\label{subsection:0.11.2}

\begin{env}[11.2.1]
\label{0.11.2.1}
Given an abelian category $\mathcal C$, we adopt the convention of
denoting by notations such as $K^\bullet$ the \emph{complexes}
$(K^i)_{i\in\mathbb Z}$ of objects of $\mathcal C$ in which the
differential is of degree $+1$, and by notations such as $K_\bullet$ the
complexes $(K_i)_{i\in\mathbb Z}$ of objects of $\mathcal C$ in which
the differential is of degree $-1$. To every complex
$K^\bullet=(K^i)$ whose differential $d$ is of degree $+1$, we
can associate a complex $K'_\bullet=(K'_i)$ by setting $K'_i=K^{-i}$,
the differential $K'_i\to K'_{i-1}$ being the operator
$d:K^{-i}\to K^{-i+1}$; and \emph{vice versa}, which will allow us, depending
on the circumstances, to consider one or the other type of complex and to
translate every result for one of the types into results for the other. Likewise
we will denote by notations such as $K^{\bullet\bullet}=(K^{ij})$ (resp.
$K_{\bullet\bullet}=(K_{ij})$) \emph{bicomplexes} of objects of
$\mathcal C$ in which \emph{both} differentials are of
degree $+1$ (resp. $-1$); again we pass from one type to the other by changing
the signs of the indices, and we have analogous notations and remarks for
arbitrary multicomplexes. The notation $K^\bullet$ or $K_\bullet$ will also be
used for \emph{graded objects} of $\mathcal C$, of type $\mathbb Z$,
which are not necessarily complexes (or which can be regarded
as such for zero differentials); for example, we will
write
$\mathrm H^\bullet(K^\bullet)=(\mathrm H^i(K^\bullet))_{i\in\mathbb Z}$ for the
\emph{cohomology} of a complex $K^\bullet$ whose differential
is of degree $+1$, and
$\mathrm H_\bullet(K_\bullet)=(\mathrm H_i(K_\bullet))_{i\in\mathbb Z}$ for
the \emph{homology} of a complex $K_\bullet$ whose differential
is of degree $-1$; when we pass from $K^\bullet$ to $K'_\bullet$ by
the operation described above, we have
$\mathrm H_i(K'_\bullet)=\mathrm H^{-i}(K^\bullet)$.

Recall in this connection that for a complex $K^\bullet$ (resp. $K_\bullet$),
we will in general write
$Z^i(K^\bullet)=\operatorname{Ker}(K^i\to K^{i+1})$ (``object of cocycles'')
and $B^i(K^\bullet)=\operatorname{Im}(K^{i-1}\to K^i)$ (``object of
coboundaries'') (resp.
$Z_i(K_\bullet)=\operatorname{Ker}(K_i\to K_{i-1})$ (``object of cycles'')
and $B_i(K_\bullet)=\operatorname{Im}(K_{i+1}\to K_i)$ (``object of
boundaries'')), so that
$\mathrm H^i(K^\bullet)=Z^i(K^\bullet)/B^i(K^\bullet)$ (resp.
$\mathrm H_i(K_\bullet)=Z_i(K_\bullet)/B_i(K_\bullet)$).

If $K^\bullet=(K^i)$ (resp. $K_\bullet=(K_i)$) is a complex in
$\mathcal C$, and $T:\mathcal C\to\mathcal C'$ a functor from $\mathcal C$
to an abelian category $\mathcal C'$, we denote by
$T(K^\bullet)$ (resp. $T(K_\bullet)$) the complex $(T(K^i))$ (resp.
$(T(K_i))$) in $\mathcal C'$.

We do not return to the definition of $\partial$-functors (T, 2.1),
except to point out that we will also say $\partial$-functor instead of
$\partial^*$-functor when the morphism $\partial$ decreases the degree by one
unit, the context having to make this point precise each time if there can be any
doubt.

Finally, we will say that a \emph{graded object} $(A_i)_{i\in\mathbb Z}$ of
$\mathcal C$ is \emph{bounded below} (resp. \emph{above})
if there exists $i_0$ such that $A_i=0$ for $i<i_0$ (resp. $i>i_0$).
\end{env}

\begin{env}[11.2.2]
\label{0.11.2.2}
Let $K^\bullet$ be a complex of $\mathcal C$ whose differential
$d$ is of degree $+1$, and suppose it equipped with a \emph{filtration}
$F(K^\bullet)=(F^p(K^\bullet))_{p\in\mathbb Z}$ consisting of graded subobjects
\oldpage[0\textsubscript{III}]{28}
of $K^\bullet$, in other words
$F^p(K^\bullet)=(K^i\cap F^p(K^\bullet))_{i\in\mathbb Z}$; moreover, we
suppose that $d(F^p(K^\bullet))\subset F^p(K^\bullet)$ for every
$p\in\mathbb Z$. Let us then recall briefly how a spectral sequence
$E(K^\bullet)$ is defined \emph{functorially} from
$K^\bullet$ (M, XV, 4 and G, I, 4.3). For $r\geq2$, the canonical morphism
$F^p(K^\bullet)/F^{p+r}(K^\bullet)\to
F^p(K^\bullet)/F^{p+1}(K^\bullet)$ defines a morphism in cohomology
\[
  \mathrm H^{p+q}(F^p(K^\bullet)/F^{p+r}(K^\bullet))
  \longrightarrow
  \mathrm H^{p+q}(F^p(K^\bullet)/F^{p+1}(K^\bullet)).
\]
We denote by $Z_r^{pq}(K^\bullet)$ the image of this morphism. Likewise, from the
exact sequence
\[
  0\to F^p(K^\bullet)/F^{p+1}(K^\bullet)
  \to F^{p-r+1}(K^\bullet)/F^{p+1}(K^\bullet)
  \to F^{p-r+1}(K^\bullet)/F^p(K^\bullet)\to0
\]
we deduce, by the cohomology exact sequence, a morphism
\[
  \mathrm H^{p+q-1}(F^{p-r+1}(K^\bullet)/F^p(K^\bullet))
  \longrightarrow
  \mathrm H^{p+q}(F^p(K^\bullet)/F^{p+1}(K^\bullet))
\]
and we denote by $B_r^{pq}(K^\bullet)$ the image of this morphism; one shows
that $B_r^{pq}(K^\bullet)\subset Z_r^{pq}(K^\bullet)$ and we take
$E_r^{pq}(K^\bullet)=Z_r^{pq}(K^\bullet)/B_r^{pq}(K^\bullet)$; we will not
spell out the definition of the $d_r^{pq}$, nor of the $\alpha_r^{pq}$.

Note here that all the $Z_r^{pq}(K^\bullet)$ and
$B_r^{pq}(K^\bullet)$, for fixed $p$, $q$, are subobjects of the same
object
$\mathrm H^{p+q}(F^p(K^\bullet)/F^{p+1}(K^\bullet))$, which is denoted
$Z_1^{pq}(K^\bullet)$; we set $B_1^{pq}(K^\bullet)=0$, so that the
preceding definitions for $Z_r^{pq}(K^\bullet)$ and
$B_r^{pq}(K^\bullet)$ also apply for $r=1$; we further set
$E_1^{pq}(K^\bullet)=Z_1^{pq}(K^\bullet)$. We further define the $d_1^{pq}$
and $\alpha_1^{pq}$ so that the conditions of
(\hyperref[0.11.1.1]{11.1.1}) are satisfied for $r=1$. On the other hand,
we define the subobjects $Z_\infty^{pq}(K^\bullet)$, the image of the morphism
\[
  \mathrm H^{p+q}(F^p(K^\bullet))\longrightarrow
  \mathrm H^{p+q}(F^p(K^\bullet)/F^{p+1}(K^\bullet))
  =E_1^{pq}(K^\bullet)
\]
and $B_\infty^{pq}(K^\bullet)$, the image of the morphism
\[
  \mathrm H^{p+q-1}(K^\bullet/F^p(K^\bullet))\longrightarrow
  \mathrm H^{p+q}(F^p(K^\bullet)/F^{p+1}(K^\bullet))
  =E_1^{pq}(K^\bullet)
\]
deduced as above from a cohomology exact sequence. We take for
$Z_\infty(E_2^{pq}(K^\bullet))$ and $B_\infty(E_2^{pq}(K^\bullet))$ the
canonical images in $E_2^{pq}(K^\bullet)$ of
$Z_\infty^{pq}(K^\bullet)$ and $B_\infty^{pq}(K^\bullet)$.

Finally, we denote by $F^p(\mathrm H^n(K^\bullet))$ the image in
$\mathrm H^n(K^\bullet)$ of the morphism
$\mathrm H^n(F^p(K^\bullet))\to\mathrm H^n(K^\bullet)$ arising from
the canonical injection $F^p(K^\bullet)\to K^\bullet$; by the cohomology exact
sequence, it is also the kernel of the morphism
$\mathrm H^n(K^\bullet)\to\mathrm H^n(K^\bullet/F^p(K^\bullet))$. We
thus define a filtration on
$E^n(K^\bullet)=\mathrm H^n(K^\bullet)$; we will not give here
the definition of the isomorphisms $\beta^{pq}$ either.
\end{env}

\begin{env}[11.2.3]
\label{0.11.2.3}
The \emph{functorial} character of $E(K^\bullet)$ is to be understood in the
following way: given two \emph{filtered} complexes $K^\bullet$,
$K^{\prime\bullet}$ of $\mathcal C$, and a morphism of complexes
$u:K^\bullet\to K^{\prime\bullet}$ \emph{compatible with the filtrations}, we
deduce from it in an evident way the morphisms $u_r^{pq}$ (for $r\geq1$) and
$u^n$, and one shows that these morphisms are compatible with the $d_r^{pq}$,
$\alpha_r^{pq}$ and $\beta^{pq}$ in the sense of
(\hyperref[0.11.1.2]{11.1.2}), and hence indeed define a morphism
$E(u):E(K^\bullet)\to E(K^{\prime\bullet})$ of spectral sequences. Moreover,
one shows that if $u$ and $v$ are morphisms
$K^\bullet\to K^{\prime\bullet}$ of the preceding type that are \emph{homotopic of order
$\leq k$}, then $u_r^{pq}=v_r^{pq}$ for $r>k$ and $u^n=v^n$ for every $n$
(M, XV, 3.1).
\end{env}

\begin{env}[11.2.4]
\label{0.11.2.4}
\oldpage[0\textsubscript{III}]{29}
Suppose that in $\mathcal C$ filtered inductive limits are
exact. Then, if the filtration $(F^p(K^\bullet))$ of $K^\bullet$ is
\emph{exhaustive}, so is the filtration
$(F^p(\mathrm H^n(K^\bullet)))$ for every $n$, since by hypothesis we have
$K^\bullet=\varinjlim_{p\to-\infty}F^p(K^\bullet)$ and the hypothesis on
$\mathcal C$ implies that cohomology commutes with inductive limits.
Moreover, for the same reason, we have
$B_\infty(E_2^{pq}(K^\bullet))=\sup_k B_k(E_2^{pq}(K^\bullet))$. We say that
the filtration $(F^p(K^\bullet))$ of $K^\bullet$ is \emph{regular} if for
every $n$ there exists an integer $u(n)$ such that
$\mathrm H^n(F^p(K^\bullet))=0$ for $p>u(n)$. This is the case in particular
when the filtration of $K^\bullet$ is \emph{discrete}. When the
filtration of $K^\bullet$ is regular and exhaustive, and filtered inductive
limits in $\mathcal C$ are exact, one shows (M, XV, 4)
that the spectral sequence $E(K^\bullet)$ is \emph{regular}.
\end{env}

\subsection{The spectral sequences of a bicomplex}
\label{subsection:0.11.3}

\begin{env}[11.3.1]
\label{0.11.3.1}
As regards the conventions concerning bicomplexes, we follow
those of (T, 2.4) rather than those of (M), the two differentials $d'$, $d''$
(of degree $+1$) of such a bicomplex
$K^{\bullet\bullet}=(K^{ij})$ being therefore assumed to \emph{commute}.
Suppose that one of the following two conditions holds:
\begin{enumerate}
  \item[$1^\circ$] Infinite direct sums exist in $\mathcal C$.
  \item[$2^\circ$] For every $n\in\mathbb Z$, there are only a
    \emph{finite} number of pairs $(p,q)$ such that $p+q=n$ and $K^{pq}\ne0$.
\end{enumerate}
Then the bicomplex $K^{\bullet\bullet}$ defines a (simple) complex
$(K^{\prime n})_{n\in\mathbb Z}$, with
$K^{\prime n}=\sum_{i+j=n}K^{ij}$, the differential $d$ (of degree
$+1$) of this complex being given by
$dx=d'x+(-1)^i d''x$ for $x\in K^{ij}$. When we speak in what follows
of the (simple) complex \emph{defined by a bicomplex $K^{\bullet\bullet}$}, it
will always be understood that one of the preceding conditions is satisfied.
We will adopt analogous conventions for multicomplexes.

We denote by $K^{i,\bullet}$ (resp. $K^{\bullet,j}$) the simple complex
$(K^{ij})_{j\in\mathbb Z}$ (resp. $(K^{ij})_{i\in\mathbb Z}$), and by
$Z_{\mathrm{II}}^p(K^{i,\bullet})$, $B_{\mathrm{II}}^p(K^{i,\bullet})$,
$\mathrm H_{\mathrm{II}}^p(K^{i,\bullet})$ (resp.
$Z_{\mathrm I}^p(K^{\bullet,j})$, $B_{\mathrm I}^p(K^{\bullet,j})$,
$\mathrm H_{\mathrm I}^p(K^{\bullet,j})$) its $p$-th objects of cocycles,
of coboundaries and of cohomology respectively; the differential
$d':K^{i,\bullet}\to K^{i+1,\bullet}$ is a morphism of complexes, which
therefore gives an operator on the cocycles, the coboundaries and the cohomology,
\[
\begin{aligned}
  d'&:Z_{\mathrm{II}}^p(K^{i,\bullet})
      \longrightarrow Z_{\mathrm{II}}^p(K^{i+1,\bullet}),\\
  d'&:B_{\mathrm{II}}^p(K^{i,\bullet})
      \longrightarrow B_{\mathrm{II}}^p(K^{i+1,\bullet}),\\
  d'&:\mathrm H_{\mathrm{II}}^p(K^{i,\bullet})
      \longrightarrow \mathrm H_{\mathrm{II}}^p(K^{i+1,\bullet}),
\end{aligned}
\]
and it is clear that for these operators,
$(Z_{\mathrm{II}}^p(K^{i,\bullet}))_{i\in\mathbb Z}$,
$(B_{\mathrm{II}}^p(K^{i,\bullet}))_{i\in\mathbb Z}$ and
$(\mathrm H_{\mathrm{II}}^p(K^{i,\bullet}))_{i\in\mathbb Z}$ are
complexes; we will denote the complex
$(\mathrm H_{\mathrm{II}}^p(K^{i,\bullet}))_{i\in\mathbb Z}$ by
$\mathrm H_{\mathrm{II}}^p(K^{\bullet\bullet})$, and its $q$-th objects of
cocycles, of coboundaries and of cohomology by
$Z_{\mathrm I}^q(\mathrm H_{\mathrm{II}}^p(K^{\bullet\bullet}))$,
$B_{\mathrm I}^q(\mathrm H_{\mathrm{II}}^p(K^{\bullet\bullet}))$ and
$\mathrm H_{\mathrm I}^q(\mathrm H_{\mathrm{II}}^p(K^{\bullet\bullet}))$.
We define similarly the complexes
$\mathrm H_{\mathrm I}^p(K^{\bullet\bullet})$ and their cohomology objects
$\mathrm H_{\mathrm{II}}^q(\mathrm H_{\mathrm I}^p(K^{\bullet\bullet}))$.
Recall on the other hand that $\mathrm H^n(K^{\bullet\bullet})$ denotes the
$n$-th cohomology object of the (simple) complex defined by
$K^{\bullet\bullet}$.
\end{env}

\begin{env}[11.3.2]
\label{0.11.3.2}
On the complex defined by a bicomplex $K^{\bullet\bullet}$, we can
consider two canonical filtrations
$(F_{\mathrm I}^p(K^{\bullet\bullet}))$ and
$(F_{\mathrm{II}}^p(K^{\bullet\bullet}))$ given by
\[
\label{0.11.3.2.1}
  F_{\mathrm I}^p(K^{\bullet\bullet})
  =\left(\sum_{\substack{i+j=n\\i\geq p}}K^{ij}\right)_{n\in\mathbb Z},
  \qquad
  F_{\mathrm{II}}^p(K^{\bullet\bullet})
  =\left(\sum_{\substack{i+j=n\\j\geq p}}K^{ij}\right)_{n\in\mathbb Z}.
  \tag{11.3.2.1}
\]
\oldpage[0\textsubscript{III}]{30}
which, by definition, are indeed graded subobjects of the (simple) complex
defined by $K^{\bullet\bullet}$, and therefore make this complex a filtered
complex; moreover, it is clear that these filtrations are exhaustive and
separated.

To each of these filtrations there corresponds a spectral sequence
(\hyperref[0.11.2.2]{11.2.2}); we will denote by
$'E(K^{\bullet\bullet})$ and $''E(K^{\bullet\bullet})$ the spectral sequences
corresponding to $(F_{\mathrm I}^p(K^{\bullet\bullet}))$ and
$(F_{\mathrm{II}}^p(K^{\bullet\bullet}))$ respectively, called the
\emph{spectral sequences of the bicomplex $K^{\bullet\bullet}$}, both
having as abutment the cohomology
$(\mathrm H^n(K^{\bullet\bullet}))$. One shows moreover (M, XV, 6) that we have
\[
\label{0.11.3.2.2}
  'E_2^{pq}(K^{\bullet\bullet})
  =\mathrm H_{\mathrm I}^p(\mathrm H_{\mathrm{II}}^q(K^{\bullet\bullet})),
  \qquad
  ''E_2^{pq}(K^{\bullet\bullet})
  =\mathrm H_{\mathrm{II}}^q(\mathrm H_{\mathrm I}^p(K^{\bullet\bullet})).
  \tag{11.3.2.2}
\]

Every morphism $u:K^{\bullet\bullet}\to K^{\prime\bullet\bullet}$ of
bicomplexes is \emph{ipso facto} compatible with the filtrations of the same
type of $K^{\bullet\bullet}$ and $K^{\prime\bullet\bullet}$, and hence defines a
morphism for each of the two spectral sequences; moreover, two
\emph{homotopic} morphisms define a homotopy \emph{of order $\leq1$} of the
corresponding filtered (simple) complexes, and hence the \emph{same} morphism for
each of the two spectral sequences (M, XV, 6.1).
\end{env}

\begin{proposition}[11.3.3]
\label{0.11.3.3}
Let $K^{\bullet\bullet}=(K^{ij})$ be a bicomplex in an abelian category
$\mathcal C$.
\begin{enumerate}
  \item[(i)] If there exist $i_0$ and $j_0$ such that $K^{ij}=0$ for $i<i_0$ or
    $j<j_0$ (resp. $i>i_0$ or $j>j_0$), the two spectral sequences
    $'E(K^{\bullet\bullet})$ and $''E(K^{\bullet\bullet})$ are
    biregular.
  \item[(ii)] If there exist $i_0$ and $i_1$ such that $K^{ij}=0$ for $i<i_0$ or
    $i>i_1$ (resp. if there exist $j_0$ and $j_1$ such that $K^{ij}=0$ for
    $j<j_0$ or $j>j_1$), the two spectral sequences
    $'E(K^{\bullet\bullet})$ and $''E(K^{\bullet\bullet})$ are
    biregular.
\end{enumerate}
Suppose moreover that in $\mathcal C$ filtered inductive limits
exist and are exact. Then:
\begin{enumerate}
  \item[(iii)] If there exists $i_0$ such that $K^{ij}=0$ for $i>i_0$ (resp. if
    there exists $j_0$ such that $K^{ij}=0$ for $j<j_0$), the sequence
    $'E(K^{\bullet\bullet})$ is regular.
  \item[(iv)] If there exists $i_0$ such that $K^{ij}=0$ for $i<i_0$ (resp. if
    there exists $j_0$ such that $K^{ij}=0$ for $j>j_0$), the sequence
    $''E(K^{\bullet\bullet})$ is regular.
\end{enumerate}
\end{proposition}

\begin{proof}
The proposition follows immediately from the definitions
(\hyperref[0.11.1.3]{11.1.3}) and from
(\hyperref[0.11.2.4]{11.2.4}), as well as from the following observations
concerning the filtration $F_{\mathrm I}$ (and the analogous observations
deduced from them for $F_{\mathrm{II}}$ by exchanging the roles of the two indices
in $K^{\bullet\bullet}$):
\begin{enumerate}
  \item[$1^\circ$] If there exists $i_0$ such that $K^{ij}=0$ for $i>i_0$, the
    filtration $F_{\mathrm I}(K^{\bullet\bullet})$ is discrete.
  \item[$2^\circ$] If there exists $i_0$ such that $K^{ij}=0$ for $i<i_0$, the
    filtration $F_{\mathrm I}(K^{\bullet\bullet})$ is co-discrete. We
    deduce immediately that the same holds for the corresponding filtration
    $F_{\mathrm I}(\mathrm H^n(K^{\bullet\bullet}))$ for every $n$;
    moreover, the definition of $B_r^{pq}$ corresponding to the filtration
    $F_{\mathrm I}(K^{\bullet\bullet})$
    (\hyperref[0.11.2.2]{11.2.2}) shows that for every pair $(p,q)$, the
    sequence $(B_r^{pq})_{r\geq2}$ is stationary.
  \item[$3^\circ$] If there exists $j_0$ such that $K^{ij}=0$ for $j<j_0$, we have
    \[
      F_{\mathrm I}^{p+r}(K^{\bullet\bullet})
      \cap\left(\sum_{i+j=n}K^{ij}\right)=0
    \]
    as soon as $p+r+j_0>n$, hence
    $Z_r^{pq}=Z_\infty(E_2^{pq})$ for $r>q-j_0+1$; on the other hand,
    $\mathrm H^n(F_{\mathrm I}^p(K^{\bullet\bullet}))=0$ for
    $p>n-j_0+1$.
  \item[$4^\circ$] If there exists $j_0$ such that $K^{ij}=0$ for $j>j_0$, we have
    \[
      F_{\mathrm I}^{p-r+1}(K^{\bullet\bullet})
      \cap\left(\sum_{i+j=n}K^{ij}\right)=\sum_{i+j=n}K^{ij}.
    \]
    \oldpage[0\textsubscript{III}]{31}
    as soon as $p-r+1+j_0<n$, hence
    $B_r^{pq}=B_\infty(E_2^{pq})$ for $r>j_0-q+1$; on the other hand,
    $\mathrm H^n(F_{\mathrm I}^p(K^{\bullet\bullet}))
    =\mathrm H^n(K^{\bullet\bullet})$ for $p+j_0<n-1$.
\end{enumerate}
\end{proof}

\begin{env}[11.3.4]
\label{0.11.3.4}
Suppose that the bicomplex $K^{\bullet\bullet}=(K^{ij})$ is such that
$K^{ij}=0$ for $i<0$ or $j<0$. We know that we can then define, for every
$p\in\mathbb Z$, a canonical ``edge homomorphism''
\[
\label{0.11.3.4.1}
  'E_2^{p0}(K^{\bullet\bullet})\longrightarrow
  \mathrm H^p(K^{\bullet\bullet})
  \tag{11.3.4.1}
\]
(M, XV, 6). Let us recall briefly that this is due, on the one hand, to the fact
that in the spectral sequence $'E(K^{\bullet\bullet})$ we have
$Z_r^{p0}=Z_{\mathrm I}^p(Z_{\mathrm{II}}^0(K^{\bullet\bullet}))$ for
$2\leq r\leq+\infty$, and on the other hand to the fact that
$\mathrm H^p(F_{\mathrm I}^{p+1}(K^{\bullet\bullet}))=0$, so that
the isomorphism
$\beta^{p0}:'E_\infty^{p0}\xrightarrow{\sim}
\mathrm H^p(F_{\mathrm I}^p)/\mathrm H^p(F_{\mathrm I}^{p+1})$ gives a
homomorphism
$'E_\infty^{p0}\to\mathrm H^p(F_{\mathrm I}^p(K^{\bullet\bullet}))
\to\mathrm H^p(K^{\bullet\bullet})$; the equality of all the $Z_r^{p0}$ then
makes it possible to define canonical homomorphisms
$'E_r^{p0}\to'E_s^{p0}$ for $r\leq s$, and in particular a homomorphism
$'E_2^{p0}\to'E_\infty^{p0}$, whence by composition the edge homomorphism
$'E_2^{p0}\to\mathrm H^p(K^{\bullet\bullet})$; moreover, one verifies immediately
that, to the class mod. $B_2^{p0}$ of an element
$z\in Z_{\mathrm{II}}^0(K^{\bullet\bullet})\subset K^{p0}$ such that $d'z=0$,
the edge homomorphism thus defined assigns in $'E_\infty^{p0}$ the
class of $z$ mod. $B_\infty^{p0}$, and then to the latter the cohomology
class of $z$ in $\mathrm H^p(K^{\bullet\bullet})$. We thus see
finally that \emph{the edge homomorphism
(\hyperref[0.11.3.4.1]{11.3.4.1}) arises, by passing to cohomology,
from the canonical injection}
$Z_{\mathrm{II}}^0(K^{\bullet\bullet})\to K^{\bullet\bullet}$ (where
$K^{\bullet\bullet}$ is considered as a simple complex). In the same way we
naturally interpret the edge homomorphism
\[
\label{0.11.3.4.2}
  ''E_2^{p0}(K^{\bullet\bullet})\longrightarrow
  \mathrm H^p(K^{\bullet\bullet})
  \tag{11.3.4.2}
\]
as arising from the canonical injection
$Z_{\mathrm I}^0(K^{\bullet\bullet})\to K^{\bullet\bullet}$.
\end{env}

\begin{env}[11.3.5]
\label{0.11.3.5}
Now let $K_{\bullet\bullet}=(K_{ij})$ be a bicomplex of
$\mathcal C$ whose two differentials are of degree $-1$. We
will then write $K_{i,\bullet}$ (resp. $K_{\bullet,j}$) for the simple complex
$(K_{ij})_{j\in\mathbb Z}$ (resp. $(K_{ij})_{i\in\mathbb Z}$),
$\mathrm H_p^{\mathrm{II}}(K_{i,\bullet})$ (resp.
$\mathrm H_p^{\mathrm I}(K_{\bullet,j})$) for its $p$-th homology object,
$\mathrm H_p^{\mathrm{II}}(K_{\bullet\bullet})$ (resp.
$\mathrm H_p^{\mathrm I}(K_{\bullet\bullet})$) for the complex
$(\mathrm H_p^{\mathrm{II}}(K_{i,\bullet}))_{i\in\mathbb Z}$ (resp.
$(\mathrm H_p^{\mathrm I}(K_{\bullet,j}))_{j\in\mathbb Z}$), and
$\mathrm H_q^{\mathrm I}(\mathrm H_p^{\mathrm{II}}(K_{\bullet\bullet}))$
(resp.
$\mathrm H_q^{\mathrm{II}}(\mathrm H_p^{\mathrm I}(K_{\bullet\bullet}))$)
for its $q$-th homology object; analogous notation for the objects of cycles
and objects of boundaries; finally, $\mathrm H_n(K_{\bullet\bullet})$ will denote
(when it exists) the $n$-th homology object of the simple complex (with
differential of degree $-1$) defined by $K_{\bullet\bullet}$.

Let $K^{\prime\bullet\bullet}=(K^{\prime ij})$, with
$K^{\prime ij}=K_{-i,-j}$, be the bicomplex with differentials of degree
$+1$ associated to $K_{\bullet\bullet}$; by definition, the
\emph{spectral sequences} of $K_{\bullet\bullet}$ are those of
$K^{\prime\bullet\bullet}$, which are written $'E(K_{\bullet\bullet})$ and
$''E(K_{\bullet\bullet})$, where however the notation is changed by
setting
\[
  'E_{pq}^r(K_{\bullet\bullet})
  ='E_r^{-p,-q}(K^{\prime\bullet\bullet}),\qquad
  ''E_{pq}^r(K_{\bullet\bullet})
  =''E_r^{-p,-q}(K^{\prime\bullet\bullet})
\]
for $2\leq r\leq\infty$. With this notation, we have
\[
  'E_{pq}^2(K_{\bullet\bullet})
  =\mathrm H_p^{\mathrm I}(\mathrm H_q^{\mathrm{II}}(K_{\bullet\bullet})),
  \qquad
  ''E_{pq}^2(K_{\bullet\bullet})
  =\mathrm H_q^{\mathrm{II}}(\mathrm H_p^{\mathrm I}(K_{\bullet\bullet})).
\]
To avoid sign errors, it will in general be preferable, for the
relations between these spectral sequences and their abutment, to go back to the
complex $K^{\prime\bullet\bullet}$. Let us note, however, the criteria
corresponding to (\hyperref[0.11.3.3]{11.3.3}):
\end{env}

\begin{env}[11.3.6]
\label{0.11.3.6}
The spectral sequences $'E(K_{\bullet\bullet})$ and
$''E(K_{\bullet\bullet})$ are \emph{biregular} in the following cases:
\begin{enumerate}
  \item[(a)] There exist $i_0$ and $j_0$ such that $K_{ij}=0$ for $i>i_0$ or
    for $j>j_0$ (resp. for $i<i_0$
    \oldpage[0\textsubscript{III}]{32}
    or for $j<j_0$).
  \item[(b)] There exist $i_0$ and $i_1$ such that $K_{ij}=0$ for $i<i_0$ and
    $i>i_1$.
  \item[(c)] There exist $j_0$ and $j_1$ such that $K_{ij}=0$ for $j<j_0$ and
    $j>j_1$.
\end{enumerate}

The sequence $'E(K_{\bullet\bullet})$ is regular if there exists $i_0$ such that
$K_{ij}=0$ for $i<i_0$, or if there exists $j_0$ such that $K_{ij}=0$ for
$j>j_0$.

The sequence $''E(K_{\bullet\bullet})$ is regular if there exists $i_0$ such that
$K_{ij}=0$ for $i>i_0$, or if there exists $j_0$ such that $K_{ij}=0$ for
$j<j_0$.
\end{env}


\subsection{Hypercohomology of a functor with respect to a complex
$K^\bullet$}
\label{subsection:0.11.4}

\begin{env}[11.4.1]
\label{0.11.4.1}
Let $\mathcal C$ be an abelian category; recall that a
\emph{right} (or \emph{cohomological}) \emph{resolution} of an object $A$ of
$\mathcal C$ is a complex of objects of $\mathcal C$, whose
differential is of degree $+1$,
\[
  0\longrightarrow L^0\longrightarrow L^1\longrightarrow L^2
  \longrightarrow\cdots
\]
equipped with a morphism $\varepsilon:A\to L^0$ called the \emph{augmentation} of the
resolution (and which can be regarded as a morphism of complexes
\[
\xymatrix{
  0\ar[r] & A\ar[r]\ar[d]_{\varepsilon} & 0\ar[r]\ar[d]
    & 0\ar[r]\ar[d] & \cdots\\
  0\ar[r] & L^0\ar[r] & L^1\ar[r] & L^2\ar[r] & \cdots
})
\]
such that the sequence
\[
  0\longrightarrow A\xrightarrow{\varepsilon}L^0\longrightarrow L^1
  \longrightarrow\cdots
\]
is exact; likewise a \emph{left} (or \emph{homological}) \emph{resolution}
of $A$ is a complex
$0\leftarrow L_0\leftarrow L_1\leftarrow\cdots$ of objects of $\mathcal C$
whose differential is of degree $-1$, equipped with an augmentation
$\varepsilon:L_0\to A$, such that the sequence
\[
  0\longleftarrow A\xleftarrow{\varepsilon}L_0\longleftarrow L_1
  \longleftarrow\cdots
\]
is exact.

When a right resolution $(L^i)_{i\geq0}$ of an object $A$ is such that
$L^i=0$ for $i\geq n+1$, we say that this resolution is of
\emph{length $\leq n$}. A left resolution of length $\leq n$ is defined
similarly. A resolution that is of length $\leq n$ for some integer
$n$ is said to be \emph{finite}.

A resolution of $A$ is said to be \emph{projective} (resp. \emph{injective}) if
the objects of $\mathcal C$ other than $A$ of which it is composed are
\emph{projective} (resp. \emph{injective}). When $\mathcal C$ is the
category of modules (left modules, for example) over a ring, we will likewise say
that a resolution of $A$ is \emph{flat} (resp. \emph{free}) when the
modules other than $A$ of which it is composed are \emph{flat} (resp.
\emph{free}).
\end{env}

\begin{env}[11.4.2]
\label{0.11.4.2}
Let $K^\bullet=(K^i)_{i\in\mathbb Z}$ be a complex of objects of $\mathcal C$
whose differential is of degree $+1$.

We call a \emph{right Cartan--Eilenberg resolution} of $K^\bullet$ the
pair consisting of a bicomplex $L^{\bullet\bullet}=(L^{ij})$ with
differentials of degree $+1$, with $L^{ij}=0$ for $j<0$, and of a morphism of
simple complexes $\varepsilon:K^\bullet\to L^{\bullet,0}$, such that the
following conditions are satisfied:
\oldpage[0\textsubscript{III}]{33}
\begin{enumerate}
  \item[(i)] For each index $i$, the sequences
    \[
    \begin{aligned}
      0\longrightarrow{}&K^i\xrightarrow{\varepsilon}L^{i0}
        \longrightarrow L^{i1}\longrightarrow\cdots\\
      0\longrightarrow{}&\mathrm B^i(K^\bullet)\xrightarrow{\varepsilon}
        \mathrm B_{\mathrm I}^i(L^{\bullet,0})
        \longrightarrow\mathrm B_{\mathrm I}^i(L^{\bullet,1})
        \longrightarrow\cdots\\
      0\longrightarrow{}&\mathrm Z^i(K^\bullet)\xrightarrow{\varepsilon}
        \mathrm Z_{\mathrm I}^i(L^{\bullet,0})
        \longrightarrow\mathrm Z_{\mathrm I}^i(L^{\bullet,1})
        \longrightarrow\cdots\\
      0\longrightarrow{}&\mathrm H^i(K^\bullet)\xrightarrow{\varepsilon}
        \mathrm H_{\mathrm I}^i(L^{\bullet,0})
        \longrightarrow\mathrm H_{\mathrm I}^i(L^{\bullet,1})
        \longrightarrow\cdots
    \end{aligned}
    \]
    are exact; in other words $(L^{i,\bullet})$,
    $(\mathrm B_{\mathrm I}^i(L^{\bullet,\bullet}))$,
    $(\mathrm Z_{\mathrm I}^i(L^{\bullet,\bullet}))$ and
    $(\mathrm H_{\mathrm I}^i(L^{\bullet,\bullet}))$ are
    \emph{resolutions} of $K^i$, $\mathrm B^i(K^\bullet)$,
    $\mathrm Z^i(K^\bullet)$ and $\mathrm H^i(K^\bullet)$ respectively.
  \item[(ii)] For each $j$, the simple complex $L^{\bullet,j}$ is
    \emph{split}, in other words, the exact sequences
    \[
      \label{0.11.4.2.1}
      0\longrightarrow\mathrm B_{\mathrm I}^i(L^{\bullet,j})
      \longrightarrow\mathrm Z_{\mathrm I}^i(L^{\bullet,j})
      \longrightarrow\mathrm H_{\mathrm I}^i(L^{\bullet,j})
      \longrightarrow0
      \tag{11.4.2.1}
    \]
    \[
      \label{0.11.4.2.2}
      0\longrightarrow\mathrm Z_{\mathrm I}^i(L^{\bullet,j})
      \longrightarrow L^{ij}
      \longrightarrow\mathrm B_{\mathrm I}^{i+1}(L^{\bullet,j})
      \longrightarrow0
      \tag{11.4.2.2}
    \]
    are split.
\end{enumerate}

One proves (M, XVII, 1.2) that if every object of $\mathcal C$ is a subobject
of an injective object, every complex $K^\bullet$ of $\mathcal C$ admits an
\emph{injective} Cartan--Eilenberg resolution, that is, one consisting of
injective objects $L^{ij}$ (condition (ii) above then implies that the
$\mathrm B_{\mathrm I}^i(L^{\bullet,j})$,
$\mathrm Z_{\mathrm I}^i(L^{\bullet,j})$ and
$\mathrm H_{\mathrm I}^i(L^{\bullet,j})$ are also injective objects).
Moreover, for every morphism $f:K^\bullet\to K'^\bullet$ of complexes of
$\mathcal C$, every Cartan--Eilenberg resolution $L^{\bullet\bullet}$ of
$K^\bullet$ and every injective Cartan--Eilenberg resolution
$L'^{\bullet\bullet}$ of $K'^\bullet$, there exists a morphism of bicomplexes
$F:L^{\bullet\bullet}\to L'^{\bullet\bullet}$ compatible with $f$ and the
augmentations, and if $f$ and $g$ are two homotopic morphisms from $K^\bullet$
to $K'^\bullet$, the corresponding morphisms from $L^{\bullet\bullet}$
to $L'^{\bullet\bullet}$ are homotopic (\emph{loc. cit.}).

When $K^\bullet$ is \emph{bounded below} (resp.
\emph{above}), we can take $L^{\bullet\bullet}$ such that
$L^{ij}=0$ for $i<i_0$ (resp. $i>i_0$) if $K^i=0$ for $i<i_0$ (resp.
$i>i_0$) (M, XVII, 1.3).

Suppose on the other hand that there exists an integer $n$ such that every object of
$\mathcal C$ admits an \emph{injective resolution of length $\leq n$};
then we may suppose that $L^{ij}=0$ for $j>n$ (M, XVII, 1.4).
\end{env}

\begin{env}[11.4.3]
\label{0.11.4.3}
Now let $T$ be a \emph{covariant additive functor} from $\mathcal C$
to an abelian category $\mathcal C'$. Given a complex
$K^\bullet$ of $\mathcal C$ and an injective Cartan--Eilenberg resolution
$L^{\bullet\bullet}$ of $K^\bullet$, suppose that the (simple) complex defined
by the bicomplex $T(L^{\bullet\bullet})$ exists
(cf. \hyperref[0.11.3.1]{11.3.1}); then the two spectral sequences
$'E(T(L^{\bullet\bullet}))$ and $''E(T(L^{\bullet\bullet}))$ of this bicomplex
are called the \emph{hypercohomology spectral sequences of $T$ with respect to the
complex $K^\bullet$}; by virtue of
(\hyperref[0.11.4.2]{11.4.2}) and
(\hyperref[0.11.3.2]{11.3.2}), they indeed depend only on
$K^\bullet$ and not on the injective Cartan--Eilenberg resolution
$L^{\bullet\bullet}$ chosen; moreover, they depend
\emph{functorially} on $K^\bullet$. They have the same abutment
$\mathrm H^\bullet(T(L^{\bullet\bullet}))$, also called
the \emph{hypercohomology of $T$ with respect to $K^\bullet$}, and denoted
$\mathrm R^\bullet T(K^\bullet)$. One shows that the $E_2$ terms of the two
preceding spectral sequences are given by
\[
  \label{0.11.4.3.1}
  'E_2^{pq}=\mathrm H^p(\mathrm R^qT(K^\bullet))
  \tag{11.4.3.1}
\]
\[
  \label{0.11.4.3.2}
  ''E_2^{pq}=\mathrm R^pT(\mathrm H^q(K^\bullet))
  \tag{11.4.3.2}
\]
\oldpage[0\textsubscript{III}]{34}
$\mathrm R^pT$ denoting as usual the $p$-th \emph{derived
functor} of $T$ for $p\in\mathbb Z$; $\mathrm R^qT(K^\bullet)$ denotes the
complex $(\mathrm R^qT(K^i))_{i\in\mathbb Z}$.

\phantomsection\egalabel{0.11.4.4}{11.4.4}
Unless expressly stated otherwise, we will henceforth suppose that every object
of $\mathcal C$ is a subobject of an injective object of $\mathcal C$, so
that injective Cartan--Eilenberg resolutions exist for every complex
of $\mathcal C$. Since $L^{ij}=0$ for $j<0$, the criteria of
(\hyperref[0.11.3.3]{11.3.3}) show that \emph{both} hypercohomology
spectral sequences of $T$ with respect to $K^\bullet$ exist and
are \emph{biregular} in each of the following two cases:
\begin{enumerate}
  \item[$1^\circ$] $K^\bullet$ is bounded below;
  \item[$2^\circ$] Every object of $\mathcal C$ admits an injective resolution
    of length at most equal to an integer $n$ (independent of the object
    under consideration).
\end{enumerate}
Indeed, in the first case we may suppose
(\hyperref[0.11.4.2]{11.4.2}) that there exists $i_0$ such that $L^{ij}=0$ for
$i<i_0$, and in the second that there exists $j_1$ such that $L^{ij}=0$ for $j>j_1$;
in each of the two cases it is moreover clear that for given $n$ there are
only finitely many pairs $(i,j)$ such that $L^{ij}\neq0$ and $i+j=n$, which
establishes our assertions.

When we suppose that in $\mathcal C'$ filtered inductive limits
exist and are exact (which implies in particular the existence in
$\mathcal C'$ of infinite direct sums), then the complex defined by the
bicomplex $T(L^{\bullet\bullet})$ exists, and the criterion
(\hyperref[0.11.3.3]{11.3.3}) shows that the sequence
$'E(T(L^{\bullet\bullet}))$ is always \emph{regular}.

When $K^\bullet$ is a complex all of whose terms $K^i$ are zero except
a single one $K^{i_0}$, $\mathrm R^nT(K^\bullet)$ is isomorphic to
$\mathrm R^{n-i_0}T(K^{i_0})$, as follows immediately from the definitions by
taking a Cartan--Eilenberg resolution $L^{\bullet\bullet}$ such that
$L^{ij}=0$ for $i\neq i_0$.

If $K^\bullet$ and $K'^\bullet$ are two complexes of $\mathcal C$, and $f$, $g$
two \emph{homotopic} morphisms from $K^\bullet$ to $K'^\bullet$, then the
morphisms $\mathrm R^\bullet T(K^\bullet)\to
\mathrm R^\bullet T(K'^\bullet)$ deduced from $f$ and $g$ are identical, and the
same holds for the morphisms of the cohomology spectral sequences.
\end{env}

\begin{proposition}[11.4.5]
\label{0.11.4.5}
Suppose that in $\mathcal C'$ filtered inductive limits exist and
are exact. If $\mathrm R^nT(K^i)=0$ for every $n>0$ and every
$i\in\mathbb Z$, we have functorial isomorphisms
\[
  \label{0.11.4.5.1}
  \mathrm R^iT(K^\bullet)\xrightarrow{\sim}
  \mathrm H^i(T(K^\bullet))\qquad(i\in\mathbb Z).
  \tag{11.4.5.1}
\]
\end{proposition}

\begin{proof}
Indeed, the only nonzero $E_2$ terms of the first spectral sequence
(\hyperref[0.11.4.3.1]{11.4.3.1}) are then
$'E_2^{p0}=\mathrm H^p(T(K^\bullet))$; in other words, this sequence is
\emph{degenerate}; since it is regular
(\hyperref[0.11.4.4]{11.4.4}), the conclusion follows from
(\hyperref[0.11.1.6]{11.1.6}).
\end{proof}

\begin{env}[11.4.6]
\label{0.11.4.6}
Now consider, for example, a \emph{covariant bifunctor}
$(M,N)\mapsto T(M,N)$ from $\mathcal C\times\mathcal C'$ to $\mathcal C''$,
$\mathcal C$, $\mathcal C'$, $\mathcal C''$ being three abelian
categories; we suppose, for simplicity, that $T$ is additive in each of its
arguments, and moreover that every object of $\mathcal C$ and every object of
$\mathcal C'$ is a subobject of an injective object, and that filtered inductive
limits exist in $\mathcal C''$ and are exact. We then define
the \emph{hypercohomology of $T$ with respect to two complexes}
$K^\bullet$, $K'^\bullet$ of $\mathcal C$ and $\mathcal C'$ respectively, with
differentials of degree $+1$, by taking for $K^\bullet$ (resp.
$K'^\bullet$) an injective Cartan--Eilenberg resolution
$L^{\bullet\bullet}$ (resp. $L'^{\bullet\bullet}$); then
$T(L^{\bullet\bullet},L'^{\bullet\bullet})$ is a quadricomplex of
$\mathcal C''$, which we regard as a \emph{bicomplex} of $\mathcal C''$
by taking as degrees of $T(L^{ij},L'^{hk})$ the integers $i+h$ and $j+k$.
The \emph{hypercohomology of $T$ with respect to $K^\bullet$ and $K'^\bullet$}
is by definition the cohomology
$\mathrm H^\bullet(T(L^{\bullet\bullet},L'^{\bullet\bullet}))$ of this
bicomplex (in other words, that of the associated simple complex), denoted
$\mathrm R^\bullet T(K^\bullet,K'^\bullet)$; it is the abutment of
two spectral sequences whose $E_2$ terms are given by
\oldpage[0\textsubscript{III}]{35}
\[
  \label{0.11.4.6.1}
  'E_2^{pq}=\mathrm H^p(\mathrm R^qT(K^\bullet,K'^\bullet))
  \tag{11.4.6.1}
\]
\[
  \label{0.11.4.6.2}
  ''E_2^{pq}=\sum_{q'+q''=q}
  \mathrm R^pT(\mathrm H^{q'}(K^\bullet),\mathrm H^{q''}(K'^\bullet))
  \tag{11.4.6.2}
\]
(cf. M, XVII, 2).

Here $\mathrm R^qT(K^\bullet,K'^\bullet)$ is the bicomplex
$(\mathrm R^qT(K^i,K'^j))_{(i,j)\in\mathbb Z\times\mathbb Z}$ and the
right-hand side of (\hyperref[0.11.4.6.1]{11.4.6.1}) is its cohomology when it is
regarded as a simple complex.

Moreover, the first spectral sequence is always regular, and the two
spectral sequences are biregular when there exists $n$ such that every object of
$\mathcal C$ and every object of $\mathcal C'$ admits an injective resolution
of length $\leq n$, or when $K^\bullet$ and $K'^\bullet$ are bounded
below; in the latter case, we can moreover omit the hypothesis that
inductive limits exist in $\mathcal C''$.

If $K_1^\bullet$, $K_1'^\bullet$ are two other complexes of $\mathcal C$ and
$\mathcal C'$ respectively, $f$, $g$ two homotopic morphisms from
$K^\bullet$ to $K_1^\bullet$, and $f'$, $g'$ two homotopic morphisms from
$K'^\bullet$ to $K_1'^\bullet$, then the morphisms
$\mathrm R^\bullet T(K^\bullet,K'^\bullet)\to
\mathrm R^\bullet T(K_1^\bullet,K_1'^\bullet)$ deduced from $f$ and $f'$ on the one
hand, and from $g$ and $g'$ on the other, are identical, and the same holds for
the morphisms of the hypercohomology spectral sequences.

This generalizes easily to an arbitrary covariant additive multifunctor.
\end{env}

\begin{proposition}[11.4.7]
\label{0.11.4.7}
Suppose that for every injective object $I$ of $\mathcal C$ (resp. $I'$ of
$\mathcal C'$), $A'\mapsto T(I,A')$ (resp. $A\mapsto T(A,I')$) is an
exact functor. Then, with the notation of
(\hyperref[0.11.4.6]{11.4.6}), we have canonical isomorphisms
\[
  \label{0.11.4.7.1}
  \mathrm R^\bullet T(K^\bullet,K'^\bullet)
  \xrightarrow{\sim}\mathrm H^\bullet(T(L^{\bullet\bullet},K'^\bullet))
  \xrightarrow{\sim}\mathrm H^\bullet(T(K^\bullet,L'^{\bullet\bullet}))
  \tag{11.4.7.1}
\]
where the last two terms are the cohomology of the simple complexes defined
by the tricomplexes $T(L^{\bullet\bullet},K'^\bullet)$ and
$T(K^\bullet,L'^{\bullet\bullet})$ respectively.
\end{proposition}

\begin{proof}
Let us define, for example, the first of these isomorphisms. The quadricomplex
$T(L^{\bullet\bullet},L'^{\bullet\bullet})$ can be regarded as a
\emph{bicomplex} by taking as degrees of $T(L^{ij},L'^{hk})$ the numbers
$i+j+h$ and $k$. Since for each $h$, $L'^{h,\bullet}$ is a
\emph{resolution} of $K'^h$, we have for this bicomplex, by virtue of the hypothesis
on $T$,
$\mathrm H_{\mathrm{II}}^q(T(L^{\bullet\bullet},L'^{\bullet\bullet}))=0$
for $q\neq0$ and
$\mathrm H_{\mathrm{II}}^0(T(L^{\bullet\bullet},L'^{\bullet\bullet}))
=T(L^{\bullet\bullet},K'^\bullet)$; the first spectral sequence of this
bicomplex is therefore \emph{degenerate}; since $L'^{hk}=0$ for $k<0$, this
sequence is moreover regular (\hyperref[0.11.3.3]{11.3.3}), and the
conclusion therefore follows from (\hyperref[0.11.1.6]{11.1.6}).
\end{proof}

We have analogous results for a covariant multifunctor in an arbitrary
number $n$ of arguments: in the computation of hypercohomology, it is not
necessary to replace \emph{all} the complexes by a Cartan--Eilenberg
resolution, but only $n-1$ of them, provided that, when we fix any
$n-1$ arguments by giving them \emph{injective} objects as values,
the covariant functor in the remaining argument is
\emph{exact}.

\subsection{Passage to the inductive limit in hypercohomology}
\label{subsection:0.11.5}

\begin{lemma}[11.5.1]
\label{0.11.5.1}
Let $K^\bullet=(K^i)_{i\in\mathbb Z}$ be a complex of $\mathcal C$, and for
every integer $r\in\mathbb Z$, let $K_{(r)}^\bullet$ be the complex such that
$K_{(r)}^i=0$ for $i<r$ and $K_{(r)}^i=K^i$ for $i\geq r$. Let $T$ be a
covariant additive functor
\oldpage[0\textsubscript{III}]{36}
from $\mathcal C$ to $\mathcal C'$ commuting with
inductive limits (we suppose that filtered inductive limits exist
and are exact in $\mathcal C$ and $\mathcal C'$). Then
$\mathrm R^\bullet T(K^\bullet)$ is isomorphic to the inductive limit
\[
  \varinjlim_{r\to-\infty}\mathrm R^\bullet T(K_{(r)}^\bullet)
\]
as $r$ tends to $-\infty$.
\end{lemma}

The construction of an injective Cartan--Eilenberg resolution of $K^\bullet$
is carried out by choosing arbitrarily, for each $i$, an injective resolution
$(X_{\mathrm B}^{ij})_{j\geq0}$ of $\mathrm B^i(K^\bullet)$ and an injective
resolution $(X_{\mathrm H}^{ij})_{j\geq0}$ of
$\mathrm H^i(K^\bullet)$; this done, the method of construction shows that
the injective resolution $(L^{ij})_{j\geq0}$ of $K^i$ and the
differentials $L^{i,j}\to L^{i+1,j}$ \emph{depend only on the resolutions}
$(X_{\mathrm B}^{i,\bullet})$, $(X_{\mathrm H}^{i,\bullet})$ and
$(X_{\mathrm B}^{i+1,\bullet})$ (M, XVII, 1.2). Now, it is clear that we have
$\mathrm B^i(K_{(r)}^\bullet)=\mathrm B^i(K^\bullet)$ and
$\mathrm H^i(K_{(r)}^\bullet)=\mathrm H^i(K^\bullet)$ for $i\geq r+1$. On
the other hand, for each $r$ we have a canonical injection
$\varphi_{r-1,r}^\bullet:K_{(r)}^\bullet\to K_{(r-1)}^\bullet$,
$\varphi_{r-1,r}^i$ being the identity for $i\neq r-1$. The preceding remark
shows that, if $L^{\bullet\bullet}=(L^{ij})$ is an injective
Cartan--Eilenberg resolution of $K^\bullet$, we can for each $r$ define an
injective Cartan--Eilenberg resolution
$L_{(r)}^{\bullet\bullet}=(L_{(r)}^{ij})$ of $K_{(r)}^\bullet$ such that
$L_{(r)}^{ij}=0$ for $i<r$ and $L_{(r)}^{ij}=L^{ij}$ for $i\geq r+1$. On the
other hand, we can define a morphism of bicomplexes
$\Phi_{r-1,r}^{\bullet\bullet}:L_{(r)}^{\bullet\bullet}\to
L_{(r-1)}^{\bullet\bullet}$ corresponding to $\varphi_{r-1,r}^\bullet$, and the
method of definition of this morphism (\emph{loc. cit.}) again shows that it
can be constructed so that $\Phi_{r-1,r}^{ij}$ is the identity for
$i\neq r$ and $i\neq r-1$. We have thus defined an inductive system
$(L_{(r)}^{\bullet\bullet})$ of bicomplexes of $\mathcal C$, of which
$L^{\bullet\bullet}$ is evidently the inductive limit as $r$ tends to
$-\infty$; since direct sums and inductive limits commute, the simple
complex associated to $L^{\bullet\bullet}$ is also the inductive limit
of the simple complex associated to $L_{(r)}^{\bullet\bullet}$. Since $T$
commutes with inductive limits by hypothesis, and so does
cohomology (because of the exactness of the functor $\varinjlim$), we indeed have
\[
  \mathrm H^\bullet(T(L^{\bullet\bullet}))=
  \varinjlim\mathrm H^\bullet(T(L_{(r)}^{\bullet\bullet}))
\]
up to isomorphism.

Lemma (\hyperref[0.11.5.1]{11.5.1}) makes it possible to extend to arbitrary complexes
$K^\bullet$, by passage to the inductive limit, properties
valid for complexes \emph{bounded below}. As a first
example, we will prove:

\begin{proposition}[11.5.2]
\label{0.11.5.2}
Under the hypotheses of (\hyperref[0.11.5.1]{11.5.1}) concerning
$\mathcal C$, $\mathcal C'$ and $T$, $\mathrm R^\bullet T(K^\bullet)$ is a
cohomological functor in the abelian category of complexes of
$\mathcal C$.
\end{proposition}

\begin{proof}
Let us show that we can reduce to the case of complexes bounded below:
if we have an exact sequence
\[
  0\longrightarrow K'^\bullet\longrightarrow K^\bullet
  \longrightarrow K''^\bullet\longrightarrow0
\]
of complexes, we evidently deduce from it for each $r$ an exact sequence
\[
  0\longrightarrow K_{(r)}'^\bullet\longrightarrow K_{(r)}^\bullet
  \longrightarrow K_{(r)}''^\bullet\longrightarrow0,
\]
whence, by hypothesis, an exact sequence
\[
  \cdots\longrightarrow\mathrm R^nT(K_{(r)}'^\bullet)
  \longrightarrow\mathrm R^nT(K_{(r)}^\bullet)
  \longrightarrow\mathrm R^nT(K_{(r)}''^\bullet)
  \xrightarrow{\partial}\mathrm R^{n+1}T(K_{(r)}'^\bullet)
  \longrightarrow\cdots
\]
these exact sequences forming an inductive system; Lemma
(\hyperref[0.11.5.1]{11.5.1}) and the exactness of the functor $\varinjlim$
prove that we have an exact sequence
\[
  \cdots\longrightarrow\mathrm R^nT(K'^\bullet)
  \longrightarrow\mathrm R^nT(K^\bullet)
  \longrightarrow\mathrm R^nT(K''^\bullet)
  \xrightarrow{\partial}\mathrm R^{n+1}T(K'^\bullet)
  \longrightarrow\cdots
\]
\end{proof}

To treat the case of complexes bounded below, we may restrict to
those such that $K^i=0$ for $i<0$; they evidently form an abelian
category $\mathcal K$.

\begin{lemma}[11.5.2.1]
\label{0.11.5.2.1}
In $\mathcal K$, let $\mathfrak J$ be the set of complexes
$Q^\bullet=(Q^i)_{i\geq0}$ having the
\oldpage[0\textsubscript{III}]{37}
following properties:
\begin{enumerate}
  \item[$1^\circ$] Every $Q^i$ is an injective object of $\mathcal C$;
  \item[$2^\circ$] For every $i\geq0$, we have
    $\mathrm Z^i(Q^\bullet)=\mathrm B^i(Q^\bullet)$, and
    $\mathrm Z^i(Q^\bullet)$ is a direct factor of $Q^i$.
\end{enumerate}
Then:
\begin{enumerate}
  \item[(i)] Every $Q^\bullet\in\mathfrak J$ is an injective object of
    $\mathcal K$.
  \item[(ii)] Every object of $\mathcal K$ is isomorphic to a subcomplex
    of a complex belonging to $\mathfrak J$.
\end{enumerate}
\end{lemma}

\begin{proof}
\textup{(i)} Let $A^\bullet=(A^i)$ be an object of $\mathcal K$,
$A'^\bullet=(A'^i)$ a subobject of $A^\bullet$, and
$Q^\bullet=(Q^i)$ an object of $\mathfrak J$, and suppose given a morphism
$f=(f^i):A'^\bullet\to Q^\bullet$, which we have to extend to a morphism
$g=(g^i):A^\bullet\to Q^\bullet$. We will use the language of the category
of modules for simplicity (cf. [27]).

Let us identify $Q^i$ with
$\mathrm B^i(Q^\bullet)\oplus\mathrm B^{i+1}(Q^\bullet)$; let us proceed by
induction on $i$, and suppose therefore that the $g^j$ have been defined for $j<i$,
compatible with the differentials $d^j:A^j\to A^{j+1}$ and
$d'^j:Q^j\to Q^{j+1}$ for $j<i-1$, and moreover such that:
\begin{enumerate}
  \item[$1^\circ$] $g^{i-1}(\mathrm Z^{i-1}(A^\bullet))\subset
    \mathrm Z^{i-1}(Q^\bullet)$;
  \item[$2^\circ$] If we set $C^j=(d^j)^{-1}(A'^{j+1})$ for every $j$, then
    $d'^{i-1}\circ g^{i-1}$ coincides with $f^i\circ d^{i-1}$ on $C^{i-1}$.
\end{enumerate}
The morphism $f^i:A'^i\to Q^i$ gives, by composition with the projections,
two morphisms
$f'^i:A'^i\to\mathrm B^i(Q^\bullet)$ and
$f''^i:A'^i\to\mathrm B^{i+1}(Q^\bullet)$. Since
$d'^{i-1}\circ g^{i-1}$ maps $A^{i-1}$ into
$\mathrm B^i(Q^\bullet)$ and vanishes on $\mathrm Z^{i-1}(A^\bullet)$, it
defines a morphism
$h^i:\mathrm B^i(A^\bullet)\to\mathrm B^i(Q^\bullet)$, and since
$d'^{i-1}\circ g^{i-1}$ coincides with $f^i\circ d^{i-1}$ on $C^{i-1}$,
$h^i$ coincides with $f'^i$ on
$\mathrm B^i(A^\bullet)\cap A'^i$. Since $\mathrm B^i(Q^\bullet)$, being a direct
factor of $Q^i$, is injective, there is a morphism
$g'^i:A^i\to\mathrm B^i(Q^\bullet)$ which coincides with $h^i$ on
$\mathrm B^i(A^\bullet)$ and with $f'^i$ on $A'^i$. Consider on the other hand
the morphism
$f'^{i+1}\circ d^i:C^i\to\mathrm B^{i+1}(Q^\bullet)$, which vanishes on
$\mathrm Z^i(A^\bullet)$; since $\mathrm B^{i+1}(Q^\bullet)$ is injective,
there is a morphism $g''^i:A^i\to\mathrm B^{i+1}(Q^\bullet)$ which coincides
with $f'^{i+1}\circ d^i$ on $C^i$ and with $0$ on
$\mathrm Z^i(A^\bullet)$. It then suffices to take $g^i=g'^i+g''^i$ in order
to be able to continue the induction.

\textup{(ii)} To embed $A^\bullet=(A^i)$ in a complex belonging to
$\mathfrak J$, we take for each $i\geq1$ an injective object $Q'^i$ of
$\mathcal C$ such that there exists an injection $f'^i:A^i\to Q'^i$. Then
set
\[
  Q^i=0\quad(i<0),\qquad Q^0=Q'^1,\qquad
  Q^i=Q'^i\oplus Q'^{i+1}\quad(i\geq1),
\]
with the evident differential. Set
$f^i=f'^i+(f'^{i+1}\circ d^i)$ for every $i$ (with $f'^i=0$ for $i\leq0$);
it is immediate that $f^i$ is injective for every $i$, and that the $f^i$ are
compatible with the differentials.
\end{proof}

\begin{corollary}[11.5.2.2]
\label{0.11.5.2.2}
Every object $K^\bullet$ of $\mathcal K$ admits a right resolution consisting
of objects of $\mathfrak J$. If $L^{\bullet\bullet}$ is such a resolution,
then for every resolution $L'^{\bullet\bullet}$ of $K^\bullet$ consisting of objects
of $\mathcal K$, there is a morphism of bicomplexes
$F:L'^{\bullet\bullet}\to L^{\bullet\bullet}$ compatible with the
augmentations, and any two such morphisms $F$, $F'$ are homotopic.
\end{corollary}

\begin{proof}
This is nothing other than (M, V, 1.1 a)) applied to the abelian category
$\mathcal K$.
\end{proof}

\begin{env}[11.5.2.3]
\label{0.11.5.2.3}
These preliminaries having been established, consider an injective
Cartan--Eilenberg resolution $L'^{\bullet\bullet}$ of $K^\bullet$ and a
resolution $L^{\bullet\bullet}$ of $K^\bullet$ consisting of objects of
$\mathfrak J$, and let us show that we have an isomorphism
\[
  \mathrm H^\bullet(T(L'^{\bullet\bullet}))\xrightarrow{\sim}
  \mathrm H^\bullet(T(L^{\bullet\bullet})).
\]
Indeed, we deduce from (\hyperref[0.11.5.2.2]{11.5.2.2}) a morphism of
bicomplexes $T(L'^{\bullet\bullet})\to T(L^{\bullet\bullet})$, and consequently
a morphism
$'E(T(L'^{\bullet\bullet}))\to{}'E(T(L^{\bullet\bullet}))$ of the first
spectral sequences of these bicomplexes. Since, by virtue of
(\hyperref[0.11.3.3]{11.3.3}), these spectral sequences are regular, it
suffices (\hyperref[0.11.1.5]{11.1.5}) to see that the preceding morphism is
an isomorphism for the $E_2$ terms, or again that
$\mathrm H_{\mathrm{II}}^q(T(L^{i,\bullet}))$ is equal to
$\mathrm R^qT(K^i)$; since $L^{i,\bullet}$ is a right resolution of
$K^i$, we are reduced to proving the
\end{env}

\oldpage[0\textsubscript{III}]{38}
\begin{lemma}[11.5.2.4]
\label{0.11.5.2.4}
If $L^\bullet=(L^i)_{i\in\mathbb Z}$ is a right resolution of an object $A$
of $\mathcal C$ such that $\mathrm R^nT(L^i)=0$ for every $i\in\mathbb Z$ and
every $n>0$, then we have
$\mathrm H^\bullet T(L^\bullet)=\mathrm R^\bullet T(A)$.
\end{lemma}

\begin{proof}
This is a particular case of (T, 2.5.2).
\end{proof}

\begin{env}[11.5.2.5]
\label{0.11.5.2.5}
The proof of (\hyperref[0.11.5.2]{11.5.2}) is now completed
immediately, since $L'^{\bullet\bullet}$ is an injective resolution of
$K^\bullet$ in the abelian category $\mathcal K$; in other words,
$K^\bullet\mapsto\mathrm H^\bullet(T(L'^{\bullet\bullet}))$ is nothing other than
the right derived cohomological functor of $T$ in the category $\mathcal K$
(T, 2.3).
\end{env}

\begin{proposition}[11.5.3]
\label{0.11.5.3}
Under the hypotheses of (\hyperref[0.11.5.1]{11.5.1}) concerning
$\mathcal C$, $\mathcal C'$ and $T$, let
$L^{\bullet\bullet}=(L^{ij})$ be a bicomplex such that $L^{ij}=0$ for $j<0$ and
such that, for every $i$, $L^{i,\bullet}$ is a resolution of $K^i$; suppose
finally that $\mathrm R^nT(L^{ij})=0$ for every pair $(i,j)$ and every $n>0$.
Then there exists a functorial isomorphism
\[
  \label{0.11.5.3.1}
  \mathrm R^\bullet T(K^\bullet)\xrightarrow{\sim}
  \mathrm H^\bullet(T(L^{\bullet\bullet})).
  \tag{11.5.3.1}
\]
\end{proposition}

\begin{proof}
Let $L_{(r)}^{\bullet\bullet}=(L_{(r)}^{ij})$ be the bicomplex such that
$L_{(r)}^{ij}=0$ for $i<r$ and $L_{(r)}^{ij}=L^{ij}$ for $i\geq r$; it is
immediate that $L^{\bullet\bullet}$ is the inductive limit of
$L_{(r)}^{\bullet\bullet}$ as $r$ tends to $-\infty$; by virtue of
the hypothesis on $T$ and of (\hyperref[0.11.5.1]{11.5.1}), it therefore suffices to
prove the proposition when $K^\bullet$ is bounded below, for
example $K^i=0$ for $i<0$, and $L^{ij}=0$ for $i<0$. Let then
$L'^{\bullet\bullet}=(L'^{ij})$ be a right resolution of $K^\bullet$ consisting
of objects of $\mathfrak J$ (\hyperref[0.11.5.2.2]{11.5.2.2}); there is a
morphism of bicomplexes
$L^{\bullet\bullet}\to L'^{\bullet\bullet}$ compatible with the
augmentations, whence a morphism
$'E(T(L^{\bullet\bullet}))\to{}'E(T(L'^{\bullet\bullet}))$ for the first
spectral sequences; Lemma (\hyperref[0.11.5.2.4]{11.5.2.4}) shows,
as in (\hyperref[0.11.5.2.3]{11.5.2.3}), that this morphism is an
\emph{isomorphism}, whence the conclusion.
\end{proof}

\begin{remark}[11.5.4]
\label{0.11.5.4}
The preceding arguments prove that the conclusions of
(\hyperref[0.11.5.2]{11.5.2}) and
(\hyperref[0.11.5.3]{11.5.3}) are valid in the category of
complexes $K^\bullet$ bounded below, \emph{without supposing that $T$
commutes with filtered inductive limits}. Moreover, when we consider
only the category $\mathcal K$ of complexes $K^\bullet$ such that
$K^i=0$ for $i<0$, the characterization of $\mathrm R^\bullet T(K^\bullet)$
as the system of \emph{right derived functors} of $T$ in $\mathcal K$
shows that this cohomological functor is \emph{universal} (T, 2.3).

Another case in which (\hyperref[0.11.5.2]{11.5.2}) is valid without making
any additional hypothesis on $T$ is the case where there exists an integer $m>0$
such that every object of $\mathcal C$ admits an injective resolution of length
$\leq m$. Indeed, in the proof of
(\hyperref[0.11.5.1]{11.5.1}), all the injective resolutions of objects
of $\mathcal C$ that occur can be taken of length $\leq m$,
from which it follows immediately that the terms of total degree $n$ of the bicomplex
$T(L_{(r)}^{\bullet\bullet})$ are equal to those of
$T(L^{\bullet\bullet})$ and finite in number, as soon as $-r$ is sufficiently large, which
implies that for every $n$,
$\mathrm H^n(T(L^{\bullet\bullet}))=
\mathrm H^n(T(L_{(r)}^{\bullet\bullet}))$ as soon as $-r$ is sufficiently large. With
the notation of (\hyperref[0.11.5.2]{11.5.2}), we therefore also have
$\mathrm R^nT(K_{(r)}^\bullet)=\mathrm R^nT(K^\bullet)$ for $-r$ sufficiently large
(depending on $n$), and likewise for $K'^\bullet$ and $K''^\bullet$, whence the
conclusion. In the same way, (\hyperref[0.11.5.3]{11.5.3}) is valid
without any additional condition on $T$ when $\mathcal C$ satisfies
the preceding hypothesis and the resolutions
$L^{i,\bullet}$ are assumed to be of length $\leq m$.
\end{remark}

\begin{env}[11.5.5]
\label{0.11.5.5}
The result of (\hyperref[0.11.5.2]{11.5.2}) generalizes to
covariant multifunctors. Consider for example the situation of
(\hyperref[0.11.4.6]{11.4.6}), where we suppose that in $\mathcal C$,
$\mathcal C'$ and $\mathcal C''$
\oldpage[0\textsubscript{III}]{39}
filtered inductive limits exist and are exact, and that $T$ commutes
with inductive limits. Then $\mathrm R^\bullet T(K^\bullet,K'^\bullet)$ is
a cohomological bifunctor in each of the complexes $K^\bullet$,
$K'^\bullet$; to see this, we reduce as in
(\hyperref[0.11.5.2]{11.5.2}) to the case where $K^\bullet$ and $K'^\bullet$ are
bounded below; then, taking for $K^\bullet$ and $K'^\bullet$
injective resolutions of the type described in
(\hyperref[0.11.5.2.2]{11.5.2.2}), we are reduced to the general property
proved in (M, V, 4.1).
\end{env}

\begin{env}[11.5.6]
\label{0.11.5.6}
Likewise, the results of (\hyperref[0.11.4.7]{11.4.7}) and
(\hyperref[0.11.5.3]{11.5.3}) generalize as follows. Suppose (under
the hypotheses of (\hyperref[0.11.5.5]{11.5.5})) that we have two
bicomplexes $L^{\bullet\bullet}=(L^{ij})$,
$L'^{\bullet\bullet}=(L'^{ij})$ such that $L^{ij}=0$ and $L'^{ij}=0$ for
$j<0$, that for every $i$, $L^{i,\bullet}$ is a resolution of $K^i$ and
$L'^{i,\bullet}$ a resolution of $K'^i$, and finally that
$\mathrm R^nT(L^{ij},L'^{hk})=0$ for $n>0$ and for every system of indices
$(i,j,h,k)$. Then we have an isomorphism, functorial in $K^\bullet$ and
$K'^\bullet$,
\[
  \label{0.11.5.6.1}
  \mathrm R^\bullet T(K^\bullet,K'^\bullet)\xrightarrow{\sim}
  \mathrm H^\bullet(T(L^{\bullet\bullet},L'^{\bullet\bullet})).
  \tag{11.5.6.1}
\]
This is established as in (\hyperref[0.11.5.3]{11.5.3}) by reducing to the
case where $K^\bullet$ and $K'^\bullet$ are bounded below.

Suppose moreover that for every pair $(i,j)$ and every pair $(h,k)$,
the functors $A\mapsto T(A,L'^{hk})$ and $A'\mapsto T(L^{ij},A')$ are
exact on $\mathcal C$ and $\mathcal C'$ respectively. Then we also have
functorial isomorphisms
\[
  \label{0.11.5.6.2}
  \mathrm R^\bullet T(K^\bullet,K'^\bullet)
  \xrightarrow{\sim}\mathrm H^\bullet(T(L^{\bullet\bullet},K'^\bullet))
  \xrightarrow{\sim}\mathrm H^\bullet(T(K^\bullet,L'^{\bullet\bullet})).
  \tag{11.5.6.2}
\]
The proof is similar to that of
(\hyperref[0.11.4.7]{11.4.7}).
\end{env}

\begin{env}[11.5.7]
\label{0.11.5.7}
Note further that the results of (\hyperref[0.11.5.5]{11.5.5}) and
(\hyperref[0.11.5.6]{11.5.6}) are valid without supposing that $T$ commutes
with inductive limits, provided that we restrict to complexes
$K^\bullet$, $K'^\bullet$ bounded below, or that we suppose that every
object of $\mathcal C$ (resp. $\mathcal C'$) admits an injective resolution of
bounded length and that in (\hyperref[0.11.5.6]{11.5.6}) the bicomplexes
$L^{\bullet\bullet}$ and $L'^{\bullet\bullet}$ have their second degree bounded
above.
\end{env}

\subsection{Hyperhomology of a functor with respect to a complex
$K_\bullet$}
\label{subsection:0.11.6}

\begin{env}[11.6.1]
\label{0.11.6.1}
Let $\mathcal C$ be an abelian category and
$K_\bullet=(K_i)_{i\in\mathbb Z}$ a complex of objects of $\mathcal C$ whose
differential is of degree $-1$. A \emph{left Cartan--Eilenberg
resolution} of $K_\bullet$ consists of a bicomplex
$L_{\bullet\bullet}=(L_{ij})$ with differentials of degree $-1$, with
$L_{ij}=0$ for $j<0$, and of a morphism of simple complexes
$\varepsilon:L_{\bullet,0}\to K_\bullet$, such that the conditions obtained
from those of (\hyperref[0.11.4.2]{11.4.2}) by ``reversing the
arrows'' are satisfied. If every object of $\mathcal C$ is a quotient of a
projective object, every complex $K_\bullet$ of $\mathcal C$ admits a
\emph{projective} Cartan--Eilenberg resolution, that is, one consisting of
projective objects $L_{ij}$, with functorial properties similar to those of
(\hyperref[0.11.4.2]{11.4.2}). Moreover, if $K_\bullet$ is bounded
below (resp. above), we may suppose that $L_{ij}=0$ for
$i<i_0$ (resp. $i>i_0$) if $K_i=0$ for $i<i_0$ (resp. $i>i_0$). If every
object of $\mathcal C$ admits a projective resolution of length $\leq n$,
we may suppose that $L_{ij}=0$ for $j>n$.
\end{env}

\begin{env}[11.6.2]
\label{0.11.6.2}
\oldpage[0\textsubscript{III}]{40}
Suppose that $T$ is a covariant additive functor from $\mathcal C$ to an
abelian category $\mathcal C'$. The definition of
\emph{the hyperhomology} $L_\bullet T(K_\bullet)$ and of the \emph{hyperhomology
spectral sequences} of $T$ with respect to a complex $K_\bullet$ of
$\mathcal C$ (when they exist) is again obtained from
(\hyperref[0.11.4.3]{11.4.3}) by ``reversing the arrows'', the $E_2$
terms of the two spectral sequences thus obtained being
\[
\label{0.11.6.2.1}
\tag{11.6.2.1}
{}'E^2_{pq}=H_p(L_qT(K_\bullet))
\]
\[
\label{0.11.6.2.2}
\tag{11.6.2.2}
{}''E^2_{pq}=L_pT(H_q(K_\bullet)),
\]
where $L_pT$ denotes the $p$-th left derived functor of $T$ for $p\geq 0$, and $0$
for $p<0$, and $L_qT(K_\bullet)$ the complex
$(L_qT(K_i))_{i\in\mathbb Z}$.

The properties of hyperhomology are not all deduced by simple
``reversal of the arrows'' from those of hypercohomology (unless
additional hypotheses of the type (AB~5*) of (T, 1.5) are made on the
category $\mathcal C'$), because of the regularity conditions of the two
preceding spectral sequences, to which the criteria of
(\hyperref[0.11.3.4]{11.3.4}) must this time be applied. These criteria show that
when we suppose that in $\mathcal C'$ filtered inductive limits
exist and are exact, then the complex defined by the bicomplex
$T(L_{\bullet\bullet})$ exists, and the second spectral sequence
$\,{}''E(T(L_{\bullet\bullet}))$ is this time regular. If we suppose
either that $K_\bullet$ is bounded below, or that there exists an integer
$n$ such that every object of $\mathcal C$ admits a projective resolution of
length $\leq n$, then the two hyperhomology spectral sequences exist
(without any hypothesis on $\mathcal C'$) and are biregular.
\end{env}

\begin{proposition}[11.6.3]
\label{0.11.6.3}
Let $\mathcal C$, $\mathcal C'$ be two abelian categories and $T$ a covariant
additive functor from $\mathcal C$ to $\mathcal C'$. Then:
\begin{enumerate}
\item[{\rm(i)}] \emph{The hyperhomology $L_\bullet T(K_\bullet)$ is a
homological functor in the abelian category of complexes of
$\mathcal C$ bounded below.}
\item[{\rm(ii)}] \emph{Let $K_\bullet$ be a complex of $\mathcal C$ bounded
below. If $L_nT(K_i)=0$ for every $n>0$ and every $i\in\mathbb Z$, we
have functorial isomorphisms}
\[
\label{0.11.6.3.1}
\tag{11.6.3.1}
L_iT(K_\bullet)\simeq H_i(T(K_\bullet))\qquad(i\in\mathbb Z).
\]
\item[{\rm(iii)}] \emph{Let $K_\bullet$ be a complex of $\mathcal C$ bounded
below. Let $L_{\bullet\bullet}=(L_{ij})$ be a bicomplex such that
$L_{ij}=0$ for $j<0$ and such that, for every $i$, $L_{i,\bullet}$ is a
resolution of $K_i$; suppose finally that $L_nT(L_{ij})=0$ for every pair
$(i,j)$ and every $n>0$. Then there exists a functorial isomorphism}
\[
\label{0.11.6.3.2}
\tag{11.6.3.2}
L_\bullet T(K_\bullet)\simeq H_\bullet(T(L_{\bullet\bullet})).
\]
\end{enumerate}
The proofs are carried out like those of
(\hyperref[0.11.5.2]{11.5.2}), (\hyperref[0.11.4.5]{11.4.5}) and
(\hyperref[0.11.5.3]{11.5.3}) in the case of complexes bounded
below. We leave the details of these arguments to the reader.
\end{proposition}

\begin{env}[11.6.4]
\label{0.11.6.4}
We have entirely analogous results for covariant multifunctors that are
additive in each of the arguments. For example, for a bifunctor $T$, we have
the two hyperhomology spectral sequences with $E_2$ terms given by
\oldpage[0\textsubscript{III}]{41}
\[
\label{0.11.6.4.1}
\tag{11.6.4.1}
{}'E^2_{pq}=H_p(L_qT(K_\bullet,K'_\bullet))
\]
\[
\label{0.11.6.4.2}
\tag{11.6.4.2}
{}''E^2_{pq}=
\sum_{q'+q''=q}L_pT(H_{q'}(K_\bullet),H_{q''}(K'_\bullet)).
\]
Here again it is the second spectral sequence that is regular, the two
sequences being biregular when complexes $K_\bullet$,
$K'_\bullet$ bounded below are concerned, or when the objects of the abelian
categories under consideration have projective resolutions of fixed
length.

Moreover:
\end{env}

\begin{proposition}[11.6.5]
\label{0.11.6.5}
Let $\mathcal C$, $\mathcal C'$, $\mathcal C''$ be three abelian
categories and $T$ a covariant biadditive bifunctor from
$\mathcal C\times\mathcal C'$ to $\mathcal C''$.
\begin{enumerate}
\item[{\rm(i)}] \emph{$L_\bullet T(K_\bullet,K'_\bullet)$ is a homological
bifunctor in each of the complexes $K_\bullet$, $K'_\bullet$ bounded below
(consisting respectively of objects of $\mathcal C$ and of objects of
$\mathcal C'$).}
\item[{\rm(ii)}] \emph{Suppose $K_\bullet$ and $K'_\bullet$ bounded
below. Let $L_{\bullet\bullet}=(L_{ij})$,
$L'_{\bullet\bullet}=(L'_{ij})$ be two bicomplexes such that $L_{ij}=0$ and
$L'_{ij}=0$ for $j<0$, such that for every $i$, $L_{i,\bullet}$ is a
resolution of $K_i$ and $L'_{i,\bullet}$ a resolution of $K'_i$, and finally such that
$L_nT(L_{ij},L'_{hk})=0$ for $n>0$ and every system $(i,j,h,k)$. Then we have
a functorial isomorphism}
\[
\label{0.11.6.5.1}
\tag{11.6.5.1}
L_\bullet T(K_\bullet,K'_\bullet)
\simeq H_\bullet(T(L_{\bullet\bullet},L'_{\bullet\bullet})).
\]
\item[{\rm(iii)}] \emph{Suppose moreover that for every pair $(i,j)$ and
every pair $(h,k)$, the functors
\[
A\longmapsto T(A,L'_{hk}),\qquad
A'\longmapsto T(L_{ij},A')
\]
are exact on $\mathcal C$ and $\mathcal C'$ respectively. Then we have
functorial isomorphisms}
\[
\label{0.11.6.5.2}
\tag{11.6.5.2}
L_\bullet T(K_\bullet,K'_\bullet)
\simeq H_\bullet(T(L_{\bullet\bullet},K'_\bullet))
\simeq H_\bullet(T(K_\bullet,L'_{\bullet\bullet})).
\]
\end{enumerate}
The proofs are analogous to those of
(\hyperref[0.11.5.5]{11.5.5}) and (\hyperref[0.11.5.6]{11.5.6}).
\end{proposition}


\subsection{Hyperhomology of a functor with respect to a bicomplex
$K_{\bullet\bullet}$}
\label{subsection:0.11.7}

\begin{env}[11.7.1]
\label{0.11.7.1}
Let $\mathcal C$ be an abelian category in which every object is a
quotient of a projective object. Consider a bicomplex
$K_{\bullet\bullet}=(K_{ij})$ formed of objects of $\mathcal C$, both of whose
degrees are bounded below; we may always restrict to the case where
$K_{ij}=0$ for $i<0$ or $j<0$, and this is what we will do from now on. We
can regard $K_{\bullet\bullet}$ as a (simple) complex formed of
objects $K_{i,\bullet}=(K_{ij})_{j\geq 0}$ of the abelian category
$\mathcal K$ of complexes of objects of $\mathcal C$ in positive degrees. It
follows from Lemma (\hyperref[0.11.5.2.1]{11.5.2.1}) (or rather from the
``dual'' lemma obtained by ``reversing the arrows'') and from (M, V, 2.2) that
$K_{\bullet\bullet}$ admits a \emph{projective Cartan--Eilenberg
resolution} in the category $\mathcal K$; such a resolution is a
tricomplex $M_{\bullet\bullet\bullet}=(M_{ijk})$ of $\mathcal C$, with
all degrees $\geq 0$, formed of projective objects, such that for every $i$,
$M_{i,\bullet\bullet}$, $\mathrm B_i^{\mathrm I}(M_{\bullet\bullet\bullet})$,
$\mathrm Z_i^{\mathrm I}(M_{\bullet\bullet\bullet})$,
$H_i^{\mathrm I}(M_{\bullet\bullet\bullet})$ constitute projective
resolutions of $K_{i,\bullet}$,
$\mathrm B_i^{\mathrm I}(K_{\bullet\bullet})$,
$\mathrm Z_i^{\mathrm I}(K_{\bullet\bullet})$,
$H_i^{\mathrm I}(K_{\bullet\bullet})$ respectively in the category
$\mathcal K$; in particular, for every pair $(i,j)$,
$M_{ij,\bullet}$ is a projective resolution of $K_{ij}$ in
$\mathcal C$.
\end{env}

\begin{proposition}[11.7.2]
\label{0.11.7.2}
Let $T$ be an additive covariant functor from $\mathcal C$ to an abelian
category $\mathcal C'$. With the notation of
(\hyperref[0.11.7.1]{11.7.1}), \emph{the homology
$H_\bullet(T(M_{\bullet\bullet\bullet}))$ of the simple complex associated to
the tricomplex $T(M_{\bullet\bullet\bullet})$ is canonically isomorphic to
the hyperhomology $L_\bullet T(K_{\bullet\bullet})$ of the simple
\oldpage[0\textsubscript{III}]{42}
complex associated to $K_{\bullet\bullet}$
(\hyperref[0.11.6.2]{11.6.2}), and is
the common abutment of six biregular spectral sequences denoted
$\mathop{}^{(t)}E$ (with $t=a,b,a',b',c$ or $d$), whose $E_2$ terms are
given by the formulas}
\[
\begin{aligned}
{}^{(a)}E^2_{pq}&=L_pT(H_q(K_{\bullet\bullet})),&
{}^{(b)}E^2_{pq}&=H_p(L_q^{\mathrm{II}}T(K_{\bullet\bullet})),\\
{}^{(a')}E^2_{pq}&=L_pT(H_q^{\mathrm I}(K_{\bullet\bullet})),&
{}^{(b')}E^2_{pq}&=H_p(L_qT(K_{\bullet\bullet})),\\
{}^{(c)}E^2_{pq}&=L_pT(H_q^{\mathrm{II}}(K_{\bullet\bullet})),&
{}^{(d)}E^2_{pq}&=H_p(L_q^{\mathrm I}T(K_{\bullet\bullet})).
\end{aligned}
\]
\end{proposition}

Recall that we use the notation $F(A_\bullet)$ to denote the
complex of the objects $F(A_i)$, for every complex $A_\bullet=(A_i)$; for
example, $L_q^{\mathrm{II}}T(K_{\bullet\bullet})$ denotes the complex
$(L_q^{\mathrm{II}}T(K_{i,\bullet}))_{i\geq0}$, where
$L_q^{\mathrm{II}}T(K_{i,\bullet})$ is the hyperhomology of index $q$ of the
functor $T$ with respect to the simple complex $K_{i,\bullet}$.

\begin{proof}
Denote by $L_\bullet$ the simple complex associated to
$K_{\bullet\bullet}$, so that
\[
L_i=\bigoplus_{r+s=i}K_{rs},
\]
and set
\[
N_{ij}=\bigoplus_{r+s=i}M_{rsj};
\]
it is clear that for each $i$, $N_{i,\bullet}$ is a projective
resolution of $L_i$ in $\mathcal C$; it therefore follows from
(\hyperref[0.11.6.3]{11.6.3}) and
(\hyperref[0.11.6.4]{11.6.4}) that we have a functorial isomorphism
\[
L_\bullet T(L_\bullet)\simeq H_\bullet(T(N_{\bullet\bullet}));
\]
since the simple complex associated to the bicomplex $T(N_{\bullet\bullet})$ is
also associated to the tricomplex $T(M_{\bullet\bullet\bullet})$, this proves
the first assertion of the statement.

Moreover, $L_\bullet T(L_\bullet)$ is the abutment of the two hyperhomology
spectral sequences (\hyperref[0.11.6.2]{11.6.2}) of $T$
relative to the simple complex $L_\bullet$, which are none other than the
sequences $\mathop{}^{(b')}E$ and $\mathop{}^{(a)}E$ respectively.

Let us now regard $M_{\bullet\bullet\bullet}$ as a bicomplex
$U_{\bullet\bullet}$ with
\[
U_{ij}=\bigoplus_{r+s=j}M_{irs};
\]
$H_\bullet(T(M_{\bullet\bullet\bullet}))$ is again the abutment of the two
spectral sequences of the bicomplex $T(U_{\bullet\bullet})$. Now, for every $i$,
$M_{i,\bullet\bullet}$ is a bicomplex satisfying the conditions of
(\hyperref[0.11.6.3]{11.6.3}) relative to the simple complex
$K_{i,\bullet}$; hence
$H_q^{\mathrm{II}}(T(U_{i,\bullet}))=
L_q^{\mathrm{II}}T(K_{i,\bullet})$, and the first spectral sequence of
$T(U_{\bullet\bullet})$ is none other than $\mathop{}^{(b)}E$. On the other
hand, for every $r$, $M_{\bullet,r,\bullet}$ is a Cartan--Eilenberg resolution
of the simple complex $K_{\bullet,r}$, the computation made in (M, XV, 2)
shows that
\[
H_q^{\mathrm I}(T(M_{\bullet,rs}))
=T(H_q^{\mathrm I}(M_{\bullet,rs})),
\]
whence
\[
H_q^{\mathrm I}(T(U_{\bullet,j}))
=\bigoplus_{r+s=j}T(H_q^{\mathrm I}(M_{\bullet,rs}));
\]
in other words, the simple complex
$H_q^{\mathrm I}(T(U_{\bullet\bullet}))$ is none other than the simple complex
associated to the bicomplex
$T(H_q^{\mathrm I}(M_{\bullet\bullet\bullet}))$. Now, for every $q$,
$H_q^{\mathrm I}(M_{\bullet\bullet\bullet})$ is a projective resolution of the
simple complex $H_q^{\mathrm I}(K_{\bullet\bullet})$ in the category
$\mathcal K$; applying (\hyperref[0.11.6.3]{11.6.3}), we see that we have
\[
H_p^{\mathrm{II}}
\bigl(T(H_q^{\mathrm I}(M_{\bullet\bullet\bullet}))\bigr)
=L_pT(H_q^{\mathrm I}(K_{\bullet\bullet})),
\]
and we thus obtain the sequence $\mathop{}^{(a')}E$. Finally, the sequences
$\mathop{}^{(c)}E$ and $\mathop{}^{(d)}E$ are obtained by interchanging the
roles of the indices $i$ and $j$ in the definition of the tricomplex
$M_{\bullet\bullet\bullet}$ and applying the preceding arguments to this new
tricomplex.
\end{proof}

We say that $L_\bullet T(K_{\bullet\bullet})$ is \emph{the hyperhomology} of
$T$ relative to the bicomplex $K_{\bullet\bullet}$.

\begin{env}[11.7.3]
\label{0.11.7.3}
\emph{Remarks}.
\begin{enumerate}
\item[{\rm(i)}] It follows from (\hyperref[0.11.6.3]{11.6.3}) that
$L_\bullet T(K_{\bullet\bullet})$ is a homological functor on the
category of bicomplexes $K_{\bullet\bullet}$ of $\mathcal C$ bounded below
in each of their degrees.
\oldpage[0\textsubscript{III}]{43}
\item[{\rm(ii)}] Let $M_{\bullet\bullet\bullet}$ be a tricomplex of
$\mathcal C$ such that for each pair $(r,s)$,
$M_{rs,\bullet}$ is a resolution of $K_{rs}$ and
$L_nT(M_{ijk})=0$ for all triples $(i,j,k)$ and every $n>0$. Then we have
an isomorphism
\[
L_\bullet T(K_{\bullet\bullet})
\simeq H_\bullet(T(M_{\bullet\bullet\bullet}));
\]
indeed, with the notation of the proof of
(\hyperref[0.11.7.2]{11.7.2}), $N_{i,\bullet}$ is a resolution of
$L_i$ such that $L_nT(N_{ij})=0$ for $n>0$ and every pair $(i,j)$, and it
suffices to apply (\hyperref[0.11.6.3]{11.6.3}, (iii)).
\item[{\rm(iii)}] The results of (\hyperref[0.11.7.2]{11.7.2}) generalize
immediately to covariant multifunctors; for example,
let $\mathcal C'$ be a second abelian category in which every object
is a quotient of a projective object, $K'_{\bullet\bullet}$ a bicomplex of
$\mathcal C'$ both of whose degrees are bounded below, and $T$ an
additive covariant bifunctor from $\mathcal C\times\mathcal C'$ to an
abelian category $\mathcal C''$. If $L_\bullet$ and $L'_\bullet$ are the
simple complexes associated to $K_{\bullet\bullet}$ and
$K'_{\bullet\bullet}$ respectively, we define the hyperhomology of $T$ with
respect to the two bicomplexes $K_{\bullet\bullet}$,
$K'_{\bullet\bullet}$ as the hyperhomology
$L_\bullet T(L_\bullet,L'_\bullet)$; applying
(\hyperref[0.11.6.4]{11.6.4}) and
(\hyperref[0.11.6.5]{11.6.5}), we have, just as in
(\hyperref[0.11.7.2]{11.7.2}), six spectral sequences abutting to this
hyperhomology, which we leave to the reader to write down.
\end{enumerate}
\end{env}

\subsection{Supplement on the cohomology of simplicial complexes}
\label{subsection:0.11.8}

\begin{env}[11.8.1]
\label{0.11.8.1}
Let $A$ be a finite set and $\Sigma(A)$ the set of finite sequences
$\sigma=(\alpha_0,\ldots,\alpha_h)$ of elements of $A$ (``simplices'' of
$A$); we set $|\sigma|=\{\alpha_0,\ldots,\alpha_h\}$; recall that the
complex of chains $\mathrm C_\bullet(A)$ is the graded free abelian group
generated by the elements of $\Sigma(A)$,
$(\alpha_0,\ldots,\alpha_h)$ being of degree $h$, with a differential
defined by
\[
d(\alpha_0,\ldots,\alpha_h)=
\sum_{j=0}^{h}(-1)^j
(\alpha_0,\ldots,\widehat{\alpha_j},\ldots,\alpha_h).
\]
The subgroup $\mathrm D_\bullet(A)$ of $\mathrm C_\bullet(A)$ generated by
the chains $\sigma=(\alpha_0,\ldots,\alpha_h)$ for which two of the
$\alpha_i$ are equal, and by the chains
$\pi(\sigma)-\varepsilon_\pi\sigma$, where, for every permutation $\pi$,
\[
\pi(\sigma)=(\alpha_{\pi(0)},\ldots,\alpha_{\pi(h)})
\]
and $\varepsilon_\pi$ is the signature of $\pi$, is a subcomplex of
$\mathrm C_\bullet(A)$ whose elements are called \emph{degenerate
chains}; we set
$\mathrm L_\bullet(A)=\mathrm C_\bullet(A)/\mathrm D_\bullet(A)$, and we have a
natural homomorphism of complexes
\[
p:\mathrm C_\bullet(A)\longrightarrow\mathrm L_\bullet(A).
\]
On the other hand, we define a homomorphism of complexes
\[
j:\mathrm L_\bullet(A)\longrightarrow\mathrm C_\bullet(A)
\]
in the following way: we totally order $A$; to the class
mod.~$\mathrm D_\bullet(A)$ of a simplex
$\sigma=(\alpha_0,\ldots,\alpha_h)$ we assign $0$ if two of the
$\alpha_i$ are equal, and $\varepsilon_\pi(\beta_0,\ldots,\beta_h)$ otherwise,
where $(\beta_0,\ldots,\beta_h)=\pi(\sigma)$ is the sequence of the $\alpha_i$
arranged in increasing order. It is clear that $p\circ
j$ is the identity of $\mathrm L_\bullet(A)$.
\end{env}

\begin{env}[11.8.2]
\label{0.11.8.2}
Let $B$ be a second finite set; if $d'$, $d''$ are the differentials of
$\mathrm C_\bullet(A)$ and $\mathrm C_\bullet(B)$, the tensor product
complex $\mathrm C_\bullet(A)\otimes\mathrm C_\bullet(B)$ can be
regarded as the free abelian group generated by the elements of
$\Sigma(A)\times\Sigma(B)$, with the differential
\[
d(\sigma,\tau)=(d'\sigma,\tau)+(-1)^{h+1}(\sigma,d''\tau)
\quad\text{if }\deg\sigma=h.
\]
The natural homomorphisms
$\mathrm C_\bullet(A)\to\mathrm L_\bullet(A)$,
$\mathrm C_\bullet(B)\to\mathrm L_\bullet(B)$ define a homomorphism
\[
p:\mathrm C_\bullet(A)\otimes\mathrm C_\bullet(B)
\longrightarrow\mathrm L_\bullet(A)\otimes\mathrm L_\bullet(B),
\]
the latter tensor product being isomorphic to
\[
\frac{\mathrm C_\bullet(A)\otimes\mathrm C_\bullet(B)}
{\mathrm C_\bullet(A)\otimes\mathrm D_\bullet(B)
+\mathrm D_\bullet(A)\otimes\mathrm C_\bullet(B)}.
\]
Similarly, by means of the homomorphisms
$\mathrm L_\bullet(A)\to\mathrm C_\bullet(A)$ and
$\mathrm L_\bullet(B)\to\mathrm C_\bullet(B)$ defined in
(\hyperref[0.11.8.1]{11.8.1}) (by means of total orders on $A$ and $B$),
we define a homomorphism
\[
j:\mathrm L_\bullet(A)\otimes\mathrm L_\bullet(B)
\longrightarrow\mathrm C_\bullet(A)\otimes\mathrm C_\bullet(B)
\]
such that $p\circ j$ is the identity.
\end{env}

\oldpage[0\textsubscript{III}]{44}
With this notation:
\begin{proposition}[11.8.3]
\label{0.11.8.3}
\emph{There exists a homotopy
\[
h:\mathrm C_\bullet(A)\otimes\mathrm C_\bullet(B)
\longrightarrow\mathrm C_\bullet(A)\otimes\mathrm C_\bullet(B),
\]
such that $h(\sigma,\tau)$ is a linear combination of pairs of simplices
$(\sigma_i,\tau_i)$ with $|\sigma_i|\subset|\sigma|$,
$|\tau_i|\subset|\tau|$, and such that, for $f=j\circ p$, we have}
\[
\label{0.11.8.3.1}
\tag{11.8.3.1}
f-1=h\circ d+d\circ h.
\]
\end{proposition}

\begin{proof}
It suffices to define $h$ on each pair $(\sigma,\tau)$ of simplices,
arguing by induction on the sum of the degrees of $\sigma$ and $\tau$, since
we can take $h=0$ when this sum is $0$. Let
\[
\omega=f(\sigma,\tau)-(\sigma,\tau)-h(d(\sigma,\tau));
\]
by virtue of the induction hypothesis and the definition of $d$, we have
\[
\omega\in\mathrm C_\bullet(|\sigma|)\otimes\mathrm C_\bullet(|\tau|).
\]
We have
\[
d\omega=f(d(\sigma,\tau))-d(\sigma,\tau)-d(h(d(\sigma,\tau)))
=h(d(d(\sigma,\tau)))=0
\]
by virtue of (\hyperref[0.11.8.3.1]{11.8.3.1}) and the induction
hypothesis. Now, we have
$H_q(\mathrm C_\bullet(A))=0$ for $q>0$ (G, I, 3.7.4), hence also
\[
H_q(\mathrm C_\bullet(A)\otimes\mathrm C_\bullet(B))=0\quad(q>0),
\]
by virtue of the Künneth formula (G, I, 5.5.2). Applying this remark with
$A$ replaced by $|\sigma|$ and $B$ by $|\tau|$, we see that there exists an
element
\[
\omega'\in\mathrm C_\bullet(|\sigma|)\otimes\mathrm C_\bullet(|\tau|)
\]
such that $\omega=d\omega'$; taking $h(\sigma,\tau)=\omega'$, we verify
(\hyperref[0.11.8.3.1]{11.8.3.1}) for the pair $(\sigma,\tau)$ and the
induction can be continued.
\end{proof}

\begin{env}[11.8.4]
\label{0.11.8.4}
The notation being that of (\hyperref[0.11.8.2]{11.8.2}), we will write
$(\sigma,\tau)\leq(\sigma',\tau')$ if $|\sigma|\subset|\sigma'|$ and
$|\tau|\subset|\tau'|$. A \emph{coefficient system} $\Gamma$ on
$\Sigma(A)\times\Sigma(B)$ consists of a family
$(\Gamma_{\sigma,\tau})$ of abelian groups, where
$\Gamma_{\sigma,\tau}$ depends only on the sets $|\sigma|$ and $|\tau|$,
and a family of homomorphisms
\[
\Gamma_{\sigma,\tau}\longrightarrow\Gamma_{\sigma',\tau'}
\quad\text{for }(\sigma,\tau)\leq(\sigma',\tau'),
\]
forming an inductive system for this preorder relation. We then define
a complex of cochains $\mathrm C^\bullet(A,B;\Gamma)$ as the set
of families $\lambda=(\lambda(\sigma,\tau))$, where $(\sigma,\tau)$ ranges over
$\Sigma(A)\times\Sigma(B)$, with
$\lambda(\sigma,\tau)\in\Gamma_{\sigma,\tau}$ for every pair
$(\sigma,\tau)$. The differential is given as follows: if
\[
d(\sigma,\tau)=\sum_i\pm(\sigma_i,\tau_i),
\]
we have $|\sigma_i|\subset|\sigma|$, $|\tau_i|\subset|\tau|$ for every $i$,
and we take
\[
d\lambda(\sigma,\tau)=\sum_i\pm\lambda_i(\sigma_i,\tau_i),
\]
where $\lambda_i(\sigma_i,\tau_i)$ denotes the canonical image of
$\lambda(\sigma_i,\tau_i)$ in $\Gamma_{\sigma,\tau}$.

We will say that a cochain $\lambda\in\mathrm C^\bullet(A,B;\Gamma)$ is
\emph{bi-alternating} if $\lambda(\sigma,\tau)=0$ whenever one of the two
simplices $\sigma$, $\tau$ has two equal terms, and if we have
\[
\lambda(\pi(\sigma),\tau)=\varepsilon_\pi\lambda(\sigma,\tau),
\qquad
\lambda(\sigma,\pi'(\tau))=\varepsilon_{\pi'}\lambda(\sigma,\tau)
\]
for arbitrary permutations $\pi$, $\pi'$ of the indices. It is clear
that these cochains generate a subcomplex
$\mathrm L^\bullet(A,B;\Gamma)$ of $\mathrm C^\bullet(A,B;\Gamma)$.
\end{env}

\begin{proposition}[11.8.5]
\label{0.11.8.5}
\emph{The canonical injection
\[
\mathrm L^\bullet(A,B;\Gamma)\longrightarrow
\mathrm C^\bullet(A,B;\Gamma)
\]
defines an isomorphism on the cohomology of these two complexes.}
\end{proposition}

\begin{proof}
Note that if $p$ and $j$ have the meaning defined in
(\hyperref[0.11.8.2]{11.8.2}), the maps
\[
{}'p:\lambda\longmapsto\lambda\circ p,\qquad
{}'j:\lambda\longmapsto\lambda\circ j
\]
are defined on $\mathrm L^\bullet(A,B;\Gamma)$ and
$\mathrm C^\bullet(A,B;\Gamma)$ respectively, the first being none other
than the canonical injection. Since $\mathop{}'j\circ{}'p$ is the identity, it
suffices to prove that $\mathop{}'p\circ{}'j$ is homotopic to the identity;
now, by (\hyperref[0.11.8.3]{11.8.3}),
$\mathop{}'h:\lambda\mapsto\lambda\circ h$ is defined on
$\mathrm C^\bullet(A,B;\Gamma)$, and we can therefore transpose the identity
(\hyperref[0.11.8.3.1]{11.8.3.1}), which gives the desired result.
\end{proof}

\begin{env}[11.8.6]
\label{0.11.8.6}
\oldpage[0\textsubscript{III}]{45}
Proposition (\hyperref[0.11.8.5]{11.8.5}) reduces the computation of the
cohomology of $\mathrm L^\bullet(A,B;\Gamma)$ to that of the cohomology of
$\mathrm C^\bullet(A,B;\Gamma)$. Recall on the other hand that the latter
is, by virtue of the Eilenberg--Zilber theorem (G, I, 3.10.2), canonically
isomorphic to the cohomology of the complex of cochains defined as follows:
we form the complex of chains $\mathrm P_\bullet(A,B)$, consisting of the
linear combinations of the $(\sigma,\tau)\in\Sigma(A)\times\Sigma(B)$ such
that $\sigma$ and $\tau$ have the same degree; the differential of this
complex is given by
\[
d(\sigma,\tau)=\sum_{j,k}(-1)^{j+k}(\sigma_j,\tau_k)
\]
if
\[
d\sigma=\sum_j(-1)^j\sigma_j,\qquad
d\tau=\sum_k(-1)^k\tau_k;
\]
we then have two canonical homomorphisms of complexes
\[
f:\mathrm P_\bullet(A,B)\longrightarrow
\mathrm C_\bullet(A)\otimes\mathrm C_\bullet(B),\qquad
g:\mathrm C_\bullet(A)\otimes\mathrm C_\bullet(B)
\longrightarrow\mathrm P_\bullet(A,B),
\]
and it is proved (\emph{loc. cit.}) that there are homotopies $h$, $h'$ such
that
\[
f\circ g-1=d\circ h+h\circ d,\qquad
g\circ f-1=d\circ h'+h'\circ d.
\]
Moreover, we have
\[
f(\sigma,\tau)\in\mathrm C_\bullet(|\sigma|)
\otimes\mathrm C_\bullet(|\tau|),\qquad
g(\sigma,\tau)\in\mathrm P_\bullet(|\sigma|,|\tau|)
\]
and the homotopies $h$, $h'$ can be chosen such that
\[
h(\sigma,\tau)\in\mathrm C_\bullet(|\sigma|)
\otimes\mathrm C_\bullet(|\tau|),\qquad
h'(\sigma,\tau)\in\mathrm P_\bullet(|\sigma|,|\tau|).
\]
This point comes from the fact that the definition of $h(\sigma,\tau)$ and
$h'(\sigma,\tau)$ can be made by induction on the sum of the degrees of
$\sigma$ and $\tau$, and from the fact that
\[
H_q(\mathrm C_\bullet(|\sigma|)\otimes\mathrm C_\bullet(|\tau|))
\quad\text{and}\quad
H_q(\mathrm P_\bullet(|\sigma|,|\tau|))
\]
are zero for $q>0$ (\emph{loc. cit.}); we then argue as in
(\hyperref[0.11.8.3]{11.8.3}) and the conclusion follows.

We then define $\mathrm P^\bullet(A,B;\Gamma)$ as the set of
families $\lambda=(\lambda(\sigma,\tau))$ where $(\sigma,\tau)$ ranges over
the pairs whose terms have the same degree, with
$\lambda(\sigma,\tau)\in\Gamma_{\sigma,\tau}$, and since we have
\[
d\sigma=\sum_j(-1)^j\sigma_j,\qquad
d\tau=\sum_k(-1)^k\tau_k,
\]
\[
d\lambda(\sigma,\tau)=
\sum_{j,k}(-1)^{j+k}\lambda(\sigma_j,\tau_k)
\]
is defined and gives the differential of the complex
$\mathrm P^\bullet(A,B;\Gamma)$. This being so, the maps
\[
{}'f:\lambda\mapsto\lambda\circ f,\quad
{}'g:\lambda\mapsto\lambda\circ g,\quad
{}'h:\lambda\mapsto\lambda\circ h,\quad
{}'h':\lambda\mapsto\lambda\circ h'
\]
are all defined by virtue of the preceding remarks;
$\mathop{}'f\circ{}'g$ and $\mathop{}'g\circ{}'f$ are therefore homotopic to
the identity, whence the desired isomorphism between the cohomology of
$\mathrm C^\bullet(A,B;\Gamma)$ and that of
$\mathrm P^\bullet(A,B;\Gamma)$.
\end{env}

\begin{env}[11.8.7]
\label{0.11.8.7}
\emph{Remark}. The same reasoning as in
(\hyperref[0.11.8.3]{11.8.3}), but applied to
$\mathrm C_\bullet(A)$ and $\mathrm L_\bullet(A)$, shows that if $j$ and $p$
are defined as in (\hyperref[0.11.8.1]{11.8.1}), $f=j\circ p$
again satisfies a relation (\hyperref[0.11.8.3.1]{11.8.3.1}), with
$|h(\sigma)|\subset|\sigma|$, from which we deduce, as in
(\hyperref[0.11.8.5]{11.8.5}), an isomorphism of the cohomology of
$\mathrm L^\bullet(A;\Gamma)$ onto that of
$\mathrm C^\bullet(A;\Gamma)$, these two complexes being defined in the
obvious way. This is the result whose proof is sketched in
(G, I, 3.8.1).
\end{env}

\begin{env}[11.8.8]
\label{0.11.8.8}
Let us now return to the notation and hypotheses of
(\hyperref[0.11.8.2]{11.8.2}), and consider a complex
$\Gamma^\bullet=(\Gamma^k)$ of coefficient systems on
$\Sigma(A)\times\Sigma(B)$: for each $(\sigma,\tau)$, the
$\Gamma^k_{\sigma,\tau}$ thus form a complex of abelian groups
$(k\in\mathbb Z)$, and we have the commutative diagrams
\[
\xymatrix{
\Gamma^k_{\sigma,\tau}\ar[r]\ar[d] &
\Gamma^{k+1}_{\sigma,\tau}\ar[d]\\
\Gamma^k_{\sigma',\tau'}\ar[r] &
\Gamma^{k+1}_{\sigma',\tau'}
}
\]
\oldpage[0\textsubscript{III}]{46}
for $(\sigma,\tau)\leq(\sigma',\tau')$. Then one verifies immediately that
\[
\mathrm C^\bullet(A,B;\Gamma^\bullet)
=\bigl(\mathrm C^h(A,B;\Gamma^k)\bigr)_{(h,k)\in\mathbb Z\times\mathbb Z}
\]
is a bicomplex of abelian groups, and
\[
\mathrm L^\bullet(A,B;\Gamma^\bullet)
=\bigl(\mathrm L^h(A,B;\Gamma^k)\bigr)
\]
is a subbicomplex of it.
\end{env}

\begin{proposition}[11.8.9]
\label{0.11.8.9}
\emph{The canonical injection
\[
\mathrm L^\bullet(A,B;\Gamma^\bullet)\longrightarrow
\mathrm C^\bullet(A,B;\Gamma^\bullet)
\]
defines an isomorphism on the cohomology of these two bicomplexes.}
\end{proposition}

\begin{proof}
Set $\mathrm C^{\bullet\bullet}=
\mathrm C^\bullet(A,B;\Gamma^\bullet)$ and
$\mathrm L^{\bullet\bullet}=
\mathrm L^\bullet(A,B;\Gamma^\bullet)$ for simplicity, and note that
since $\mathrm C^{hk}=\mathrm L^{hk}=0$ for $h<0$, the second spectral
sequences of these bicomplexes are regular
(\hyperref[0.11.3.3]{11.3.3}); the homomorphism
$\mathrm L^{\bullet\bullet}\to\mathrm C^{\bullet\bullet}$ therefore gives a
morphism of spectral sequences
\[
{}''E(\mathrm L^{\bullet\bullet})\longrightarrow
{}''E(\mathrm C^{\bullet\bullet})
\]
which, for the $E_2$ terms, reduces to
\[
\label{0.11.8.9.1}
\tag{11.8.9.1}
H^p_{\mathrm{II}}(H^q_{\mathrm I}(\mathrm L^{\bullet\bullet}))
\longrightarrow
H^p_{\mathrm{II}}(H^q_{\mathrm I}(\mathrm C^{\bullet\bullet})).
\]
But for every $k\in\mathbb Z$, it follows from
(\hyperref[0.11.8.5]{11.8.5}) that the homomorphism
$H^q_{\mathrm I}(\mathrm L^{\bullet,k})\to
H^q_{\mathrm I}(\mathrm C^{\bullet,k})$ is bijective; the conclusion
therefore follows from (\hyperref[0.11.1.5]{11.1.5}).
\end{proof}

\begin{env}[11.8.10]
\label{0.11.8.10}
Similarly, with the notation of (\hyperref[0.11.8.6]{11.8.6}), we have
canonical homomorphisms of bicomplexes
\[
\mathrm C^\bullet(A,B;\Gamma^\bullet)\longrightarrow
\mathrm P^\bullet(A,B;\Gamma^\bullet)
\]
(with the obvious notation), and the same reasoning as in
(\hyperref[0.11.8.9]{11.8.9}), based this time on
(\hyperref[0.11.8.6]{11.8.6}), shows that this homomorphism again gives
an isomorphism in cohomology.
\end{env}

\subsection{A lemma on complexes of finite type}
\label{subsection:0.11.9}

\begin{proposition}[11.9.1]
\label{0.11.9.1}
Let $\mathcal C$ be an abelian category, and $K'$ and $K''$ subsets of
the set of objects of $\mathcal C$ such that $K''\subset K'$ and
satisfying the following conditions:
\begin{enumerate}
\item[{\rm(i)}] \emph{For every object $A'\in K'$ and every epimorphism
$u:A\to A'$ in $\mathcal C$, there exist an object $B\in K''$ and a morphism
$v:B\to A$ such that $u\circ v$ is an epimorphism.}
\item[{\rm(ii)}] \emph{For every pair of objects $A\in K'$, $B\in K'$ and
every epimorphism $u:A\to B$, $\operatorname{Ker}(u)$ belongs to $K'$.}
\item[{\rm(iii)}] \emph{The product of two objects of $K''$ belongs to
$K''$.}
\end{enumerate}
\emph{Let $P_\bullet=(P_i)_{i\in\mathbb Z}$ be a complex in $\mathcal C$
such that $H_i(P_\bullet)\in K'$ for every $i$ and such that there exists $d$
with $H_i(P_\bullet)=0$ for $i<d$. Then there exist a complex
$Q_\bullet=(Q_i)_{i\in\mathbb Z}$ in $\mathcal C$ such that $Q_i\in K''$
for every $i$ and $Q_i=0$ for $i<d$, and a morphism
$u:Q_\bullet\to P_\bullet$ of complexes such that the corresponding morphism
$H_\bullet(Q_\bullet)\to H_\bullet(P_\bullet)$ is an isomorphism.}
\end{proposition}

\begin{proof}
Let us first prove the following consequence of property (i):

\emph{(i bis) Let $u:C\to B$ be an epimorphism in $\mathcal C$, $A$ an
object of $K'$, and $v:A\to B$ a morphism in $\mathcal C$; then there exist an
object $D\in K''$, an epimorphism $u':D\to A$ and a morphism $v':D\to C$
such that the diagram}
\[
\xymatrix{
D\ar[r]^{u'}\ar[d]_{v'} & A\ar[d]^{v}\\
C\ar[r]_{u} & B
}
\tag{11.9.1.1}\label{0.11.9.1.1}
\]
\emph{is commutative.}

Indeed, consider the fibre product $C\times_B A$ in $\mathcal C$ and the
canonical projections $p:C\times_B A\to C$, $q:C\times_B A\to A$, which make
the diagram
\oldpage[0\textsubscript{III}]{47}
\[
\xymatrix{
C\times_B A\ar[r]^{q}\ar[d]_{p} & A\ar[d]^{v}\\
C\ar[r]_{u} & B
}
\tag{11.9.1.2}\label{0.11.9.1.2}
\]
commutative. We know ([27], p.~1-12) that the cokernel of $q$ is the
quotient of $A$ by $v^{-1}(u(C))$; since $u$ is an epimorphism, $u(C)=B$ and
$v^{-1}(u(C))=A$, so $q$ is an epimorphism; it then suffices to apply
(i) to the epimorphism $q:C\times_B A\to A$: there exist an object
$D\in K''$ and a morphism $w:D\to C\times_B A$ such that $q\circ w$ is an
epimorphism; we take $u'=q\circ w$, $v'=p\circ w$.

This being so, to prove the proposition, we proceed by induction.
Suppose that, for some $i\geq d-1$, the objects $Q_j$, the morphisms
$d_j:Q_j\to Q_{j-1}$ and the morphisms $u_j:Q_j\to P_j$ have been
constructed for $j\leq i$, in such a way that $Q_j=0$ for $j<d$,
$d_{j-1}\circ d_j=0$ and
$d_j\circ u_j=u_{j-1}\circ d_j$ for $j\leq i$; moreover, we suppose that
the following conditions are satisfied:
\begin{enumerate}
\item[$(\mathrm I_i)$] We have $Q_j\in K''$ for $j\leq i$ and
$B_j(Q_\bullet)\in K'$ for $j<i$.
\item[$(\mathrm{II}_i)$] For $j<i$, the homomorphism
$H_j(Q_\bullet)\to H_j(P_\bullet)$ deduced from the family $(u_k)_{k\leq i}$
is an isomorphism.
\item[$(\mathrm{III}_i)$] The composite morphism
\[
v_i:Z_i(Q_\bullet)\longrightarrow Z_i(P_\bullet)
\longrightarrow H_i(P_\bullet)
\]
(where the left arrow is the restriction of $u_i$ and the right arrow
the canonical morphism) is an epimorphism.
\end{enumerate}
Note that, by (ii), $Z_i(Q_\bullet)$, the kernel of the epimorphism
$Q_i\to B_{i-1}(Q_\bullet)$, belongs to $K'$ by virtue of hypothesis
$(\mathrm I_i)$. We further deduce from (ii) that
$N_i=\operatorname{Ker}(v_i)$ also belongs to $K'$, taking into account
hypothesis $(\mathrm{III}_i)$. By virtue of (i bis), there exist a
$Q'_{i+1}\in K''$, an epimorphism $d'_{i+1}:Q'_{i+1}\to N_i$ and a
morphism $u'_{i+1}:Q'_{i+1}\to P_{i+1}$ such that the diagram
\[
\xymatrix{
Q'_{i+1}\ar[r]^{d'_{i+1}}\ar[d]_{u'_{i+1}} &
N_i\ar[d]^{u_i}\\
P_{i+1}\ar[r]_{d_{i+1}} & B_i(P_\bullet)
}
\tag{11.9.1.3}\label{0.11.9.1.3}
\]
is commutative.

Since the canonical morphism $Z_{i+1}(P_\bullet)\to H_{i+1}(P_\bullet)$ is
an epimorphism and $H_{i+1}(P_\bullet)\in K'$ by hypothesis, it follows
from (i) that there exist an object $Q''_{i+1}\in K''$ and a morphism
$u''_{i+1}:Q''_{i+1}\to Z_{i+1}(P_\bullet)$ such that the composite
$Q''_{i+1}\to Z_{i+1}(P_\bullet)\to H_{i+1}(P_\bullet)$ is an
epimorphism. If we take $d''_{i+1}:Q''_{i+1}\to N_i$ equal to $0$, the
diagram
\[
\xymatrix{
Q''_{i+1}\ar[r]^{d''_{i+1}}\ar[d]_{u''_{i+1}} &
N_i\ar[d]^{u_i}\\
Z_{i+1}(P_\bullet)\ar[r]_{d_{i+1}} & P_i
}
\tag{11.9.1.4}\label{0.11.9.1.4}
\]
is commutative, the lower horizontal arrow being $0$.
\oldpage[0\textsubscript{III}]{48}
Then take
\[
Q_{i+1}=Q'_{i+1}\times Q''_{i+1},\qquad
d_{i+1}=d'_{i+1}+d''_{i+1},\qquad
u_{i+1}=u'_{i+1}+u''_{i+1}.
\]
Since $d_{i+1}(Q_{i+1})=d'_{i+1}(Q'_{i+1})=N_i\subset Z_i(Q_\bullet)$, we have
$d_i\circ d_{i+1}=0$ and, with the usual notation,
$B_i(Q_\bullet)=N_i$, which verifies $(\mathrm I_{i+1})$. The commutativity
of the diagrams (\hyperref[0.11.9.1.3]{11.9.1.3}) and
(\hyperref[0.11.9.1.4]{11.9.1.4}) shows that
$d_{i+1}\circ u_{i+1}=u_i\circ d_{i+1}$. By definition of $N_i$, the
morphism
\[
H_i(Q_\bullet)=Z_i(Q_\bullet)/N_i\longrightarrow H_i(P_\bullet)
\]
deduced from the system of the $u_k$ ($k\leq i+1$) is the morphism deduced
from $v_i$ by passage to the quotients, so it is an isomorphism since $v_i$
is an epimorphism, whence $(\mathrm{II}_{i+1})$. Finally, we have
$Q''_{i+1}\subset Z_{i+1}(Q_\bullet)$ by definition; the choice of
$u''_{i+1}$ shows that the morphism
\[
v_{i+1}:Z_{i+1}(Q_\bullet)\longrightarrow Z_{i+1}(P_\bullet)
\longrightarrow H_{i+1}(P_\bullet)
\]
is an epimorphism, since its restriction to $Q''_{i+1}$ already is, whence
$(\mathrm{III}_{i+1})$. The induction can therefore be continued, and the
proposition is proved.
\end{proof}

\begin{corollary}[11.9.2]
\label{0.11.9.2}
Let $A$ be a Noetherian ring (not necessarily commutative) and
$P_\bullet=(P_i)_{i\in\mathbb Z}$ a complex of right $A$-modules. Suppose
that the $H_i(P_\bullet)$ are $A$-modules of finite type and that there
exists $d$ such that $H_i(P_\bullet)=0$ for $i<d$. Then there exist a
complex $Q_\bullet=(Q_i)_{i\in\mathbb Z}$ formed of free right $A$-modules of
finite rank, such that $Q_i=0$ for $i<d$, and a homomorphism
$u:Q_\bullet\to P_\bullet$ of complexes such that the homomorphism
$H_\bullet(Q_\bullet)\to H_\bullet(P_\bullet)$ corresponding to $u$ is
bijective.
\end{corollary}

\begin{proof}
We apply (\hyperref[0.11.9.1]{11.9.1}), taking for $\mathcal C$ the
category of right $A$-modules, and for $K'$ (resp. $K''$) the set of
$A$-modules of finite type (resp. the set of free $A$-modules of finite
rank); the verification of conditions (i), (ii) and (iii) of
(\hyperref[0.11.9.1]{11.9.1}) is immediate, taking into account the
hypothesis that $A$ is Noetherian.
\end{proof}

\begin{env}[11.9.3]
\label{0.11.9.3}
\emph{Remarks}.
\begin{enumerate}
\item[{\rm(i)}] Under the conditions of
(\hyperref[0.11.9.2]{11.9.2}), suppose moreover that the $P_i$ are flat
right $A$-modules. Then, for every left $A$-module $M$,
the homomorphism of complexes
\[
u\otimes 1:Q_\bullet\otimes_A M\longrightarrow P_\bullet\otimes_A M
\]
again defines an isomorphism
$H_\bullet(Q_\bullet\otimes_A M)\simeq
H_\bullet(P_\bullet\otimes_A M)$ of the homology, as we will see in
Chapter~III.
\item[{\rm(ii)}] The conclusion of (\hyperref[0.11.9.2]{11.9.2}) is no
longer necessarily true when $A$ is not assumed Noetherian; indeed,
applying it to a complex reduced to $0$ except for a single term, we would
conclude that every right $A$-module of finite type admits a resolution
by free modules of finite type, which is not true in general
(cf. Bourbaki, \emph{Alg. comm.}, chap.~I, \S~2, exerc.~6).

However, instead of assuming $A$ Noetherian, we can assume only that
the $H_i(P_\bullet)$ have a finite $\infty$-presentation (cf. Chapter~IV).
\end{enumerate}
\end{env}

\subsection{Euler--Poincaré characteristic of a complex of modules of
finite length}
\label{subsection:0.11.10}

\begin{env}[11.10.1]
\label{0.11.10.1}
Let $A$ be a ring (not necessarily commutative) and
\[
\label{0.11.10.1.1}
\tag{11.10.1.1}
M^\bullet:0\longrightarrow M^0\longrightarrow M^1
\longrightarrow\cdots\longrightarrow M^n\longrightarrow0
\]
a complex of left $A$-modules of finite length. We call the
\emph{Euler--Poincaré characteristic} of this complex the number
\oldpage[0\textsubscript{III}]{49}
\[
\label{0.11.10.1.2}
\tag{11.10.1.2}
\chi(M^\bullet)=\sum_{i=0}^{n}(-1)^i\operatorname{long}(M^i).
\]
\end{env}

\begin{proposition}[11.10.2]
\label{0.11.10.2}
\emph{For every finite complex $M^\bullet$ of left $A$-modules of finite
length, we have
\[
\chi(M^\bullet)=\chi(H^\bullet(M^\bullet))
\]
($H^\bullet(M^\bullet)$ being regarded as a complex with the trivial
differential). In particular, if the sequence
(\hyperref[0.11.10.1.1]{11.10.1.1}) is exact, we have
$\chi(M^\bullet)=0$.}
\end{proposition}

\begin{proof}
For brevity, set $B^i=B^i(M^\bullet)$, $Z^i=Z^i(M^\bullet)$,
$H^i=H^i(M^\bullet)=Z^i/B^i$; the $B^i$, $Z^i$, $H^i$ are of finite
length. From the exact sequences
\[
0\longrightarrow B^i\longrightarrow Z^i\longrightarrow H^i
\longrightarrow0,
\qquad
0\longrightarrow Z^i\longrightarrow M^i\longrightarrow B^{i+1}
\longrightarrow0
\]
we obtain the relations
\[
\operatorname{long}(Z^i)=\operatorname{long}(H^i)+
\operatorname{long}(B^i),
\qquad
\operatorname{long}(M^i)=\operatorname{long}(Z^i)+
\operatorname{long}(B^{i+1}),
\]
whence
\[
\operatorname{long}(M^i)-\operatorname{long}(H^i)=
\operatorname{long}(B^{i+1})+\operatorname{long}(B^i).
\]
Multiplying this relation by $(-1)^i$ and summing the relations
obtained for $0\leq i\leq n$, noting that
$B^0=B^{n+1}=0$, we obtain the desired equality.
\end{proof}

\begin{corollary}[11.10.3]
\label{0.11.10.3}
Let $E=(E_r^{pq})$ be a spectral sequence in the category of modules over a
ring $A$. Suppose that the $E_2^{pq}$ are $A$-modules of finite
length and that there are only finitely many pairs $(p,q)$ such that
$E_2^{pq}\neq0$. Then the Euler--Poincaré characteristics
$\chi(E_r^\bullet)$ of all the complexes
$E_r^\bullet=(E_r^{(n)})_{n\in\mathbb Z}$
(\hyperref[0.11.1.1]{11.1.1}) are all equal. If, moreover, the sequence
$E$ is weakly convergent and if we set
\[
E_\infty^{(n)}=\bigoplus_{p+q=n}E_\infty^{pq}\quad(n\in\mathbb Z),
\]
we also have $\chi(E_\infty^\bullet)=\chi(E_2^\bullet)$,
$E_\infty^\bullet=(E_\infty^{(n)})_{n\in\mathbb Z}$ being regarded as
a complex with trivial differential.
\end{corollary}

\begin{proof}
Note first that if $E_2^{pq}=0$, then $E_r^{pq}=0$ for
$2\leq r<+\infty$, so all the complexes $E_r^\bullet$ are finite and
formed of $A$-modules of finite length; the relation
$\chi(E_r^\bullet)=\chi(E_{r+1}^\bullet)$ for every finite $r$ therefore
follows from (\hyperref[0.11.10.2]{11.10.2}) and from the isomorphism
between $H^\bullet(E_r^\bullet)$ and $E_{r+1}^\bullet$ (as complexes with
trivial differential). The hypothesis that the $E_r^{pq}$ are of finite length
implies that for every pair $(p,q)$, the sequences
$(B_k(E_2^{pq}))_{k\geq2}$ and $(Z_k(E_2^{pq}))_{k\geq2}$ are stationary;
the hypothesis that $E$ is weakly convergent and that $E_2^{pq}=0$ except
for finitely many pairs $(p,q)$ therefore implies that there exists an
integer $r\geq2$ such that $E_\infty^{pq}=E_r^{pq}$ for every pair $(p,q)$;
whence the assertion concerning $\chi(E_\infty^\bullet)$.
\end{proof}

